\documentclass[../../main.tex]{subfiles}
\graphicspath{{\subfix{../../image/}}}

\begin{document}

\section{Sobolev空间}\label{section:ch05-Sobolev空间}

\subsection{弱导数}\label{subsection:ch05-弱导数}

\begin{definition}
设 \(X\) 为拓扑空间，\(A,B\subseteq X\)。若
\[
\overline{A}\text{ 是 }X\text{ 中的紧集，且 }\overline{A}\subseteq B,
\]
则称 \(A\) \textbf{紧包含于} \(B\)，记作
\[
A\subset\subset B
\quad(\text{或 } A\Subset B).
\]
特别地,当 $B$ 是开集时，$A \subset\subset B$ 等价于 $\overline{A}$ 是紧集，且 $\mathrm{dist}(\overline{A}, \partial B) > 0$。
\end{definition}
\begin{remark}
引入“紧包含”这一概念，核心动机在于将“全局问题”转化为“局部问题”，并为各种分析工具提供合法的操作空间。
具体原因如下:

(i)磨光的几何限制。正如之前所分析的，磨光算子的定义要求函数在点 $x$ 的完整 $\varepsilon$-邻域内有定义，且由\hyperref[theorem:appC-光滑算子的性质]{光滑算子的性质}知,磨光后的支集会向外膨胀 $\varepsilon$的厚度，即 $\mathrm{spt}(u^\varepsilon) \subseteq \mathrm{spt}(u) + \overline{B(0,\varepsilon)}$。如果我们在 $V \subset\subset U$ 上工作，并选取磨光半径 $0 < \varepsilon < \mathrm{dist}(V, \partial U)$，那么 $V$ 的 $\varepsilon$-邻域将完全落在 $U$ 的内部。这就避开了“零延拓”带来的问题：若在边界附近直接做零延拓，卷积会在边界附近导数爆炸，导致磨光序列在 $W^{k,p}$ 范数下无法收敛，且会破坏原有的边界条件。如果只在紧包含的子集上磨光，这个子集的边界仍在定义域内部，那么磨光就是合法的.例如\hyperref[theorem:实值-Kirszbraun扩张定理]{扩张定理}的证明.

(ii)截断函数与分部积分的需求。在能量方法中，我们经常需要分部积分来导出先验估计。然而，分部积分会产生边界项，而我们对函数在边界上的行为往往缺乏控制。通过选取嵌套的紧包含 $V \subset\subset W \subset\subset U$，我们可以构造一个截断函数 $\zeta \in C_c^\infty(U)$，使得 $\zeta \equiv 1$ 于 $V$，而 $\zeta \equiv 0$ 于 $U \setminus W$。于是，在对乘以 $\zeta$ 后的函数进行分部积分时，边界项会因 $\zeta$ 在边界附近为零而彻底消失。这使得我们可以在局部区域 $V$ 上自由地使用微积分基本定理和能量估计，而不必担心边界带来的麻烦。

(iii)定义局部 Sobolev 空间与推导内部正则性。许多偏微分方程的解在区域内部具有很好的光滑性，但在边界上可能由于边界条件或区域几何形状而表现不佳。空间 $W_{\mathrm{loc}}^{k,p}(U)$ 的定义是：对任意 $V \subset\subset U$，都有 $u \in W^{k,p}(V)$。这里的紧包含将“整体”切分成了“局部”，使得我们可以逐块地证明内部正则性定理（如内部 $H^2$ 估计），再用有限覆盖定理将局部结果拼凑成全局结论。

综上所述，引入紧包含为了在边界行为难以控制的情况下，通过紧包含的子集，确保磨光、截断、分部积分能够合法。
\end{remark}

\begin{proposition}\label{proposition:紧集的相对紧开邻域}
设 $A\subseteq\mathbb{R}^n$ 为开集, $B\subseteq A$ 为紧集. 则存在开集 $U\subseteq\mathbb{R}^n$, 使得
\begin{align*}
B\subseteq U\subseteq\overline{U}\subseteq A,
\end{align*}
且 $\overline{U}$ 为紧集. 换言之, 存在包含 $B$ 的开集 $U$, 使得$U\subset \subset A$.
\end{proposition}
\begin{proof}

若 $B=\varnothing$, 取 $U=\varnothing$ 即可. 以下设 $B\neq\varnothing$.

若 $A=\mathbb{R}^n$, 取 $r=1$. 若 $A\neq\mathbb{R}^n$, 则由\subcref{Real Analysis-proposition:点集间的距离的性质}{Real Analysis-proposition:点集间的距离的性质-3}知
\begin{align*}
\delta\triangleq\operatorname{dist}(B,\partial A)>0,
\end{align*}
取 $r=\frac{\delta}{2}$,定义
\begin{align*}
d_B(x)\triangleq\operatorname{dist}(x,B)
=\inf\limits_{b\in B}\|x-b\|,\quad
U\triangleq\{x\in\mathbb{R}^n:d_B(x)<r\}.
\end{align*}
由三角不等式, 对任意 $x,y\in\mathbb{R}^n$, 均有
\begin{align*}
|d_B(x)-d_B(y)|\leqslant\|x-y\|.
\end{align*}
因此$d_B$连续, 从而 $U$ 为开集. 又因为 $d_B(b)=0$ 对每个 $b\in B$ 成立, 所以 $B\subseteq U$. 由$d_B$的连续性还有
\begin{align*}
\overline{U}\subseteq\{x\in\mathbb{R}^n:d_B(x)\leqslant r\}.
\end{align*}

下面证明 $\overline{U}\subseteq A$. 当 $A=\mathbb{R}^n$ 时显然成立. 当 $A\neq\mathbb{R}^n$ 时, 任取 $x\in\overline{U}$. 由 $B$ 的紧性, 存在 $b\in B$, 使得
\begin{align*}
\|x-b\|=d_B(x)\leqslant r.
\end{align*}
假设 $x\notin A$, 考虑线段参数化
\begin{align*}
\gamma(t)\triangleq b+t(x-b),\quad t\in[0,1].
\end{align*}
集合
\begin{align*}
T\triangleq\{t\in[0,1]:\gamma(t)\notin A\}
\end{align*}
是 $[0,1]$ 中的非空闭集, 因而存在最小元$t_0\triangleq \mathrm{inf}\,T\in T$. 由于 $\gamma(0)=b\in A$ 且 $A$ 为开集, 有 $t_0>0$. 根据 $t_0$ 的定义, 当 $0\leqslant t<t_0$ 时, $\gamma(t)\in A$, 而 $\gamma(t_0)\notin A$. 由 $\gamma$ 的连续性和\subcref{Real Analysis-theorem:边界和开集基本性质}{Real Analysis-theorem:边界和开集基本性质-5}可知
\begin{align*}
\gamma(t_0)\in\overline{A}\setminus A=\partial A.
\end{align*}
于是
\begin{align*}
\delta
=\operatorname{dist}(B,\partial A)
\leqslant\|\gamma(t_0)-b\|=t_0\|x-b\|
\leqslant r
<\delta,
\end{align*}
矛盾!因此 $x\in A$, 从而 $\overline{U}\subseteq A$.

最后证明 $\overline{U}$ 紧. 由于 $B$ 紧, 故 $B$ 有界, 可取 $R>0$, 使得 $\|b\|\leqslant R$ 对所有 $b\in B$ 成立. 对任意 $x\in\overline{U}$, 仍由 $B$ 的紧性选取 $b\in B$, 使得 $\|x-b\|=d_B(x)\leqslant r$, 则
\begin{align*}
\|x\|\leqslant\|x-b\|+\|b\|\leqslant r+R,
\end{align*}
故 $\overline{U}$ 有界. 又 $\overline{U}$ 为闭集, 由 \hyperref[Real Analysis-theorem:Heine - Borel有限子覆盖定理]{Heine--Borel定理}可知 $\overline{U}$ 为紧集.
\end{proof}

\begin{definition}
设 $U \subseteq \mathbb{R}^n$ 是一个非空开集,记 $C^{\infty}_{c}(U)$ 为所有支集含于 $U$ 中的无穷可微函数 $\phi:U\to\mathbb{R}$ 所构成的空间。我们有时把属于 $C^{\infty}_{c}(U)$ 的函数 $\phi$ 称为\textbf{检验函数}。
\end{definition}

\begin{definition}\label{definition:local_integrable_2026}
设 $U \subseteq \mathbb{R}^n$ 是一个非空开集,$1\leqslant p\leqslant \infty$,定义
\begin{align*}
L_{\mathrm{loc}}^{p}\left( U \right) \triangleq \left\{ u:U\rightarrow \mathbb{R} \mid \text{对每个}V\subset \subset U,u\in L^p\left( V \right) \right\} .
\end{align*}
称在$L_{\mathrm{loc}}^{p}$中的函数为\textbf{局部可积函数}或\textbf{局部可和函数}.

在 $L_{\mathrm{loc}}^p(U)$ 中定义等价关系 $\sim$ 如下:
\begin{align*}
f \sim g\iff f(x)=g(x)\quad \text{a.e.}\,x \in U.
\end{align*}
将商空间$L_{\mathrm{loc}}^p(U)/\sim$仍记作 $L_{\mathrm{loc}}^p(U)$. 该商空间中的元素为等价类 $[f]$, 线性运算定义为 $[f]+[g] \triangleq [f+g]$, $\lambda [f] \triangleq [\lambda f]$, 这些运算是良定义的. 因此 $L_{\mathrm{loc}}^p(U)$也是一个线性空间.
\end{definition}
\begin{remark}
$L_{\mathrm{loc}}^p(U)$和$L_{\mathrm{loc}}^p(U)/\sim$都用$L_{\mathrm{loc}}^p(U)$表示虽然是符号滥用，但是一般并不会造成歧义.今后，我们可以将$L_{\mathrm{loc}}^p(U)$中几乎处处相等的函数视为同一个元素, 即 $L_{\mathrm{loc}}^p(U)$ 是由等价类构成的线性空间.
\end{remark}

\begin{proposition}[整体可积与局部可积的关系]\label{proposition:整体可积与局部可积的关系}
设 $U\subseteq\mathbb{R}^n$ 为可测集, $1\leqslant q\leqslant p\leqslant\infty$. 则
\begin{enumerate}[(1)]
\item\label{proposition:整体可积与局部可积的关系-1} \(L^p(U)\subseteq L^p_{\mathrm{loc}}(U)
\subseteq L^q_{\mathrm{loc}}(U)\).

\item 存在可测集 \(U\) 及 \(f\in L^p_{\mathrm{loc}}(U)\)，使得 \(f\notin L^p(U)\)。
\end{enumerate}
\end{proposition}
\begin{remark}
局部可积或局部本质有界只要求在 \(U\) 的每个紧子集上有相应的有限性或本质有界性，不控制无穷远处的衰减与边界附近的奇性。因此局部可积或局部本质有界一般不能推出整体可积或整体本质有界。特别地，当 \(U\) 本身为紧集时，\(L^p_{\mathrm{loc}}(U)=L^p(U)\)（\(1\leqslant p\leqslant\infty\)），此时反向包含对紧集 \(U\) 成立.
\end{remark}
\begin{proof}

先证 (1)。分两种情形。
当 \(1\leqslant p<\infty\) 时，任取 \(f\in L^p(U)\)，则
\begin{align*}
\int_U |f(x)|^p\,\mathrm{d}x<+\infty.
\end{align*}
设 \(K\subseteq U\) 为任意紧集。由于 \(K\subseteq U\)，有
\begin{align*}
\int_K |f(x)|^p\,\mathrm{d}x
\leqslant \int_U |f(x)|^p\,\mathrm{d}x
<+\infty,
\end{align*}
故 \(f\in L^p(K)\)，从而 \(f\in L^p_{\mathrm{loc}}(U)\)。因此
\begin{align*}
L^p(U)\subseteq L^p_{\mathrm{loc}}(U).
\end{align*}

当 \(p=\infty\) 时，任取 \(f\in L^\infty(U)\)，则存在 \(M\geqslant0\)，使得
\begin{align*}
|f(x)|\leqslant M
\end{align*}
对几乎处处的 \(x\in U\) 成立。设 \(K\subseteq U\) 为任意紧集。限制在 \(K\) 上仍有
\begin{align*}
|f(x)|\leqslant M
\end{align*}
对几乎处处的 \(x\in K\) 成立，故 \(f\in L^\infty(K)\)，从而 \(f\in L^\infty_{\mathrm{loc}}(U)\)。因此
\begin{align*}
L^\infty(U)\subseteq L^\infty_{\mathrm{loc}}(U).
\end{align*}

再证 (2)。反向不成立的反例见下面的例子。

取 \(U\triangleq\mathbb{R}\)，\(f(x)\triangleq x\)。对任意紧集 \(K\subseteq\mathbb{R}\)，\(K\) 有界，故存在 \(M_K>0\)，使得对一切 \(x\in K\)，有
\begin{align*}
|f(x)|=|x|\leqslant M_K.
\end{align*}
若 \(1\leqslant p<\infty\)，则
\begin{align*}
\int_K |f(x)|^p\,\mathrm{d}x
\leqslant M_K^p\,m(K)<+\infty,
\end{align*}
故 \(f\in L^p_{\mathrm{loc}}(\mathbb{R})\)。若 \(p=\infty\)，则 \(|f|\leqslant M_K\) 在 \(K\) 上成立，故 \(f\in L^\infty_{\mathrm{loc}}(\mathbb{R})\)。因此对一切 \(1\leqslant p\leqslant\infty\)，都有 \(f\in L^p_{\mathrm{loc}}(\mathbb{R})\)。

但是，若 \(1\leqslant p<\infty\)，则
\begin{align*}
\int_{\mathbb{R}} |f(x)|^p\,\mathrm{d}x
=\int_{\mathbb{R}} |x|^p\,\mathrm{d}x
=+\infty,
\end{align*}
故 \(f\notin L^p(\mathbb{R})\)。若 \(p=\infty\)，则 \(f\) 在 \(\mathbb{R}\) 上无界，故 \(f\notin L^\infty(\mathbb{R})\)。因此
\begin{align*}
L^p_{\mathrm{loc}}(\mathbb{R})\nsubseteq L^p(\mathbb{R})
\end{align*}
对一切 \(1\leqslant p\leqslant\infty\) 成立。

下证$L^p_{\mathrm{loc}}(U)\subseteq L^q_{\mathrm{loc}}(U).$
任取 $f\in L^p_{\mathrm{loc}}(U)$, 并任取紧集 $K\subseteq U$. 则 $f\in L^p(K)$, 且由于 $K$ 有界, 有 $m(K)<\infty$.从而由\subcref{Real Analysis-proposition:Lp空间包含关系}{Real Analysis-proposition:Lp空间包含关系-1}知$f\in L^q(K)$,故$f\in L^q_{\mathrm{loc}}(U).$
\end{proof}

\begin{definition}\label{definition:ch05-弱偏导数}
设 $U \subseteq \mathbb{R}^n$ 是一个非空开集, $u,v\in L^{1}_{\mathrm{loc}}(U)$，$\alpha$ 为一多重指标。我们称 $v\in L^{1}_{\mathrm{loc}}(U)$ 为 $u$ 的 \textbf{$\alpha$ 阶弱偏导数}，记作
$D^{\alpha}u=v,$
只要对一切检验函数 $\phi\in C^{\infty}_{c}(U)$ 都有
\begin{align}
\int_{U}uD^{\alpha}\phi\,\mathrm{d}x=(-1)^{|\alpha|}\int_{U}v\phi\,\mathrm{d}x.
\label{align:ch05-弱导数定义}
\end{align}
\end{definition}
\begin{remark}
换言之，给定 $u$ 后，若恰好存在一个函数 $v$ 对一切 $\phi$ 都满足式~\eqref{align:ch05-弱导数定义}，我们就称在弱意义下 $D^{\alpha}u=v$。若不存在这样的函数 $v$，则 $u$ 不具有弱 $\alpha$ 阶偏导数。
\end{remark}
\begin{remark}
\textbf{弱导数定义的动机:} 设给定函数 $u\in C^{1}(U)$。则若 $\phi\in C^{\infty}_{c}(U)$，由分部积分公式可见
\begin{align}
\int_{U}u\phi_{x_{i}}\,\mathrm{d}x=-\int_{U}u_{x_{i}}\phi\,\mathrm{d}x \qquad (i=1,\cdots,n).
\label{align:ch05-partbyparts1}
\end{align}
由于 $\phi$ 在 $U$ 中具有紧支集，因而在 $\partial U$ 附近为零，故没有边界项。更一般地，若 $k$ 为正整数，$u\in C^{k}(U)$，且 $\alpha=(\alpha_{1},\cdots,\alpha_{n})$ 为满足 $|\alpha|=\alpha_{1}+\cdots+\alpha_{n}=k$ 的多重指标，则
\begin{align}
\int_{U}uD^{\alpha}\phi\,\mathrm{d}x=(-1)^{|\alpha|}\int_{U}D^{\alpha}u\,\phi\,\mathrm{d}x.
\label{align:ch05-partbyparts2}
\end{align}
该等式成立是因为
\begin{align*}
D^{\alpha}\phi=\frac{\partial^{\alpha_{1}}}{\partial x_{1}^{\alpha_{1}}}\cdots\frac{\partial^{\alpha_{n}}}{\partial x_{n}^{\alpha_{n}}}\phi,
\end{align*}
我们可将式~\eqref{align:ch05-partbyparts1} 应用 $|\alpha|$ 次。

下面考察对 $u\in C^{k}(U)$ 成立的式~\eqref{align:ch05-partbyparts2}，并追问：即使 $u$ 并非 $k$ 次连续可微，它的某种变形是否仍可能成立。只要 $u$ 局部可和，\eqref{align:ch05-partbyparts2}式的左端就有意义；问题在于若 $u$ 不是 $C^{k}$，则\eqref{align:ch05-partbyparts2}式右端的“$D^{\alpha}u$”一词就没有显然的含义。我们通过追问是否存在一个局部可和函数 $v$，使式~\eqref{align:ch05-partbyparts2} 仍成立而以 $v$ 代替 $D^{\alpha}u$，来解决这一困难.
\end{remark}

\begin{lemma}[变分法基本原理]\label{lemma:变分法基本原理}
设 $U\subseteq\mathbb{R}^n$ 是非空开集, $w\in L_{\mathrm{loc}}^1(U)$. 若对任意函数 $\phi\in C_c^\infty(U)$, 都有
\begin{align*}
\int_U w(x)\phi(x)\,\mathrm{d}x=0,
\end{align*}
则 $w(x)=0$ 对几乎处处的 $x\in U$ 成立.
\end{lemma}
\begin{proof}

反设 $w$ 在 $U$ 上不是几乎处处为零. 由于
\begin{align*}
\{x\in U:w(x)\neq 0\}
=
\{x\in U:w(x)>0\}\cup\{x\in U:w(x)<0\},
\end{align*}
正值集合与负值集合中至少有一个具有正测度. 必要时将 $w$ 换成 $-w$, 这不改变题设的积分恒等式. 因而不妨设
\begin{align*}
E\triangleq\{x\in U:w(x)>0\},
\quad |E|>0.
\end{align*}
对每个正整数 $m$, 定义
\begin{align*}
E_m\triangleq
\left\{x\in E:w(x)\geqslant\frac{1}{m}\right\}.
\end{align*}
由于
\begin{align*}
E=\bigcup\limits_{m=1}^{\infty}E_m,
\end{align*}
若所有 $E_m$ 都是零测集, 则 $E$ 也是零测集, 与 $|E|>0$ 矛盾. 因此存在正整数 $m_0$, 使得 $|E_{m_0}|>0$.

由\hyperref[Real Analysis-theorem:Lebesgue测度的内正则性]{Lebesgue测度的内正则性}, 可以选取紧集 $K\subseteq E_{m_0}$, 使得
$|K|>0.$
令 $c\triangleq\frac{1}{m_0}$, 则
\begin{align}
w(x)\geqslant c>0,
\quad x\in K.
\label{eq:variational-positive-on-compact-7b2f}
\end{align}
又由\cref{proposition:紧集的相对紧开邻域}, 存在开集 $W$, 使得
\begin{align*}
K\subseteq W\subseteq\overline{W}\subseteq U,
\end{align*}
且 $\overline{W}$ 为紧集. 由于 $w\in L_{\mathrm{loc}}^1(U)$, 有
\begin{align*}
\int_W|w(x)|\,\mathrm{d}x
\leqslant
\int_{\overline{W}}|w(x)|\,\mathrm{d}x
<\infty.
\end{align*}
由 \hyperref[Real Analysis-theorem:积分的绝对连续性]{Lebesgue 积分的绝对连续性}, 存在 $\delta>0$, 使得对任意可测集 $A\subseteq W$, 只要 $|A|<\delta$, 就有
\begin{align}
\int_A|w(x)|\,\mathrm{d}x
<\frac{c|K|}{2}.
\label{eq:variational-absolute-continuity-4e9a}
\end{align}

由于 $K$ 可测且具有有限测度, 由 \hyperref[Real Analysis-theorem:定理2.14]{Lebesgue 测度的外正则性}, 存在开集 $O\supseteq K$, 使得
\begin{align*}
|O\setminus K|<\delta.
\end{align*}
令 $G\triangleq O\cap W$. 则 $G$ 是包含 $K$ 的开集, 且
\begin{align*}
|G\setminus K|
\leqslant |O\setminus K|
<\delta.
\end{align*}
由于 $G\setminus K\subseteq W$, 由 \eqref{eq:variational-absolute-continuity-4e9a} 得
\begin{align}
\int_{G\setminus K}|w(x)|\,\mathrm{d}x
<\frac{c|K|}{2}.
\label{eq:variational-small-error-8c3d}
\end{align}

现在构造支撑包含于 $G$ 的光滑截断函数. 再次应用\cref{proposition:紧集的相对紧开邻域}, 选取开集 $V$, 使得
\begin{align*}
K\subseteq V\subseteq\overline{V}\subseteq G,
\end{align*}
且 $\overline{V}$ 为紧集. 由紧集与不相交闭集之间的正距离性质,
\begin{align*}
d_1\triangleq
\operatorname{dist}(K,\mathbb{R}^n\setminus V)>0,\quad
d_2\triangleq
\operatorname{dist}(\overline{V},\mathbb{R}^n\setminus G)>0.
\end{align*}
取 $0<\varepsilon<\min\{d_1,d_2\}$. 设 $\eta$ 为标准磨光核, 即
\begin{align*}
\eta\in C_c^\infty(B(0,1)),
\quad \eta\geqslant 0,
\quad \int_{\mathbb{R}^n}\eta(x)\,\mathrm{d}x=1,
\end{align*}
并令
\begin{align*}
\eta_\varepsilon(x)
\triangleq
\frac{1}{\varepsilon^n}
\eta\left(\frac{x}{\varepsilon}\right),\quad
\phi\triangleq\eta_\varepsilon*\chi_V,
\end{align*}
其中 $\chi_V$ 视为定义在整个 $\mathbb{R}^n$ 上的特征函数.

由\hyperref[theorem:appC-光滑算子的性质]{光滑算子的性质}, 有
\begin{align}
0\leqslant\phi(x)\leqslant 1,
\quad x\in\mathbb{R}^n.
\label{eq:variational-cutoff-bound-2f6b}
\end{align}
当 $x\in K$ 且 $y\in\operatorname{spt}\eta_\varepsilon$ 时, 有 $|y|\leqslant\varepsilon<d_1$, 因而 $x-y\in V$. 所以
\begin{align}
\phi(x)
=
\int_{\mathbb{R}^n}
\eta_\varepsilon(y)\chi_V(x-y)\,\mathrm{d}y=
\int_{\mathbb{R}^n}\eta_\varepsilon(y)\,\mathrm{d}y=1,
\quad x\in K.\label{eq:variational-cutoff-one-5a1c}
\end{align}
此外, 由 $\varepsilon<d_2$ 和\hyperref[theorem:appC-光滑算子的性质]{光滑算子的性质\ref{theorem:appC-光滑算子的性质-6}}得
\begin{align*}
\operatorname{spt}\phi\subseteq \mathrm{spt}\,\chi _V+\overline{B\left( 0,\varepsilon \right) }
\subseteq
\left\{
x\in\mathbb{R}^n:
\operatorname{dist}(x,\overline{V})\leqslant\varepsilon
\right\}\subseteq G.
\end{align*}
由于 $\overline{V}$ 为紧集, 且距离函数连续, 集合
$\left\{
x\in\mathbb{R}^n:
\operatorname{dist}(x,\overline{V})\leqslant\varepsilon
\right\}$
为紧集. 因此
\begin{align}
\phi\in C_c^\infty(G).
\label{eq:variational-cutoff-support-9d4e}
\end{align}
从而$\phi=0$ 于 $U\setminus G$, 从而
\begin{align}
\int_U w(x)\phi(x)\,\mathrm{d}x
=
\int_G w(x)\phi(x)\,\mathrm{d}x.
\label{eq:variational-integral-restriction-6c8a}
\end{align}
又由于 $K\subseteq G$, 且
\begin{align*}
G=K\cup(G\setminus K),
\quad K\cap(G\setminus K)=\varnothing,
\end{align*}
因此由\hyperref[Real Analysis-theorem:积分对定义域的可数可加性]{积分对可测集合的可加性}, 再结合\eqref{eq:variational-cutoff-one-5a1c}式得
\begin{align}
\int_G w(x)\phi(x)\,\mathrm{d}x
=
\int_K w(x)\phi(x)\,\mathrm{d}x
+\int_{G\setminus K}w(x)\phi(x)\,\mathrm{d}x
=
\int_K w(x)\,\mathrm{d}x
+\int_{G\setminus K}w(x)\phi(x)\,\mathrm{d}x,
\label{eq:variational-integral-decomposition-3b7f}
\end{align}
综合\eqref{eq:variational-positive-on-compact-7b2f}\eqref{eq:variational-small-error-8c3d}\eqref{eq:variational-cutoff-bound-2f6b}\eqref{eq:variational-integral-restriction-6c8a}\eqref{eq:variational-integral-decomposition-3b7f}式, 得
\begin{align*}
\int_U w(x)\phi(x)\,\mathrm{d}x
=
\int_K w(x)\,\mathrm{d}x
+\int_{G\setminus K}w(x)\phi(x)\,\mathrm{d}x
\geqslant
c|K|-\int_{G\setminus K}|w(x)|\,\mathrm{d}x
>
c|K|-\frac{c|K|}{2}
=\frac{c|K|}{2}>0.
\end{align*}
这与题设对任意 $\phi\in C_c^\infty(U)$ 均有积分为零矛盾!因此 $w=0$ 在 $U$ 上几乎处处成立.
\end{proof}

\begin{lemma}[弱导数的唯一性]\label{lemma:ch05-弱导数唯一性}
设 $U \subseteq \mathbb{R}^n$ 是一个非空开集,若$u\in L_{\mathrm{loc}}^1(U)$ 的 $\alpha$ 阶弱偏导数若存在，则在相差一个零测集的意义下唯一。
\end{lemma}
\begin{proof}
设 $v,\widetilde{v}\in L^{1}_{\mathrm{loc}}(U)$ 满足对一切 $\phi\in C^{\infty}_{c}(U)$，
\begin{align*}
\int_{U}uD^{\alpha}\phi\,\mathrm{d}x=(-1)^{|\alpha|}\int_{U}v\phi\,\mathrm{d}x=(-1)^{|\alpha|}\int_{U}\widetilde{v}\phi\,\mathrm{d}x.
\end{align*}
于是对一切 $\phi\in C^{\infty}_{c}(U)$，
\begin{align*}
\int_{U}(v-\widetilde{v})\phi\,\mathrm{d}x=0,
\end{align*}
因而由\hyperref[lemma:变分法基本原理]{变分法基本原理}知$v-\widetilde{v}=0$ 几乎处处成立。
\end{proof}

\begin{example}\label{example:ch05-弱导例1}
取 $n=1$，$U=(0,2)$，令
\begin{align*}
u(x)=
\begin{cases}
x, & 0<x\leqslant 1,\\
1, & 1\leqslant x<2.
\end{cases}
\end{align*}
求$u$的一阶弱导数.
\end{example}
\begin{remark}
这里例子说明:有尖点的函数通常可以有弱导数.
\end{remark}
\begin{solution}
定义
\begin{align*}
v(x)=
\begin{cases}
1, & 0<x\leqslant 1,\\
0, & 1<x<2.
\end{cases}
\end{align*}
我们证明在弱意义下 $u'=v$。为此任取 $\phi\in C^{\infty}_{c}(U)$。需证
\begin{align*}
\int_{0}^{2}u\phi'\,\mathrm{d}x=-\int_{0}^{2}v\phi\,\mathrm{d}x.
\end{align*}
按弱导数定义，由$u$的连续性，分部积分计算得
\begin{align*}
&\int_{0}^{2}u\phi'\,\mathrm{d}x
=\int_{0}^{1}x\phi'\,\mathrm{d}x+\int_{1}^{2}\phi'\,\mathrm{d}x
=-\int_{0}^{1}\phi\,\mathrm{d}x+\phi(1)-\phi(1)
=-\int_{0}^{2}v\phi\,\mathrm{d}x
\\
&\Longleftrightarrow \int_0^1{\left( v-1 \right) \phi \,\mathrm{d}x}+\int_1^2{v\phi \,\mathrm{d}x}=0，
\end{align*}
因此$v(x)$就是$u$的一阶弱导数.
\end{solution}

\begin{example}\label{example:ch05-弱导例2}
取 $n=1$，$U=(0,2)$，令
\begin{align*}
u(x)=
\begin{cases}
x, & 0<x\leqslant 1,\\
2, & 1<x<2.
\end{cases}
\end{align*}
求$u$的弱导数,若不存在,请证明.
\end{example}
\begin{remark}
这里例子说明:有跳跃间断点的函数通常没有有弱导数.
\end{remark}
\begin{solution}

我们证明 $u$ 不存在由 $L_{\mathrm{loc}}^1(U)$ 函数表示的弱导数.
反设存在 $v\in L_{\mathrm{loc}}^1(U)$, 使得对任意 $\phi\in C_c^\infty(U)$, 都有
\begin{align}
\int_0^2u(x)\phi'(x)\,\mathrm{d}x
=-\int_0^2v(x)\phi(x)\,\mathrm{d}x.
\label{align:ch05-弱导不存在}
\end{align}
由于 $\phi$ 在 $0$ 和 $2$ 的邻域内恒为零, 分部积分可得
\begin{align*}
\int_0^2u(x)\phi'(x)\,\mathrm{d}x
=\int_0^1x\phi'(x)\,\mathrm{d}x
+2\int_1^2\phi'(x)\,\mathrm{d}x
=\phi(1)-\int_0^1\phi(x)\,\mathrm{d}x-2\phi(1)=-\int_0^1\phi(x)\,\mathrm{d}x-\phi(1).
\end{align*}
结合 \eqref{align:ch05-弱导不存在}式, 得
\begin{align}
\phi(1)
=\int_0^2v(x)\phi(x)\,\mathrm{d}x
-\int_0^1\phi(x)\,\mathrm{d}x,
\qquad \phi\in C_c^\infty(U).
\label{align:ch05-弱导例2-点值恒等式}
\end{align}

下面构造测试函数制造矛盾. 取标准磨光子 $\eta\in C_c^\infty(\mathbb{R})$, 满足
\begin{align*}
\eta\geqslant0,\qquad
\operatorname{spt}\eta\subseteq[-1,1],\qquad
\int_{\mathbb{R}}\eta(x)\,\mathrm{d}x=1,
\end{align*}
并对每个 $\varepsilon>0$ 定义
\begin{align*}
\eta_\varepsilon(x)
\triangleq\frac{1}{\varepsilon}\eta\left(\frac{x}{\varepsilon}\right).
\end{align*}
于是
\begin{align*}
\eta_\varepsilon\geqslant0,\qquad
\operatorname{spt}\eta_\varepsilon\subseteq[-\varepsilon,\varepsilon],
\qquad
\int_{\mathbb{R}}\eta_\varepsilon(x)\,\mathrm{d}x=1.
\end{align*}

对 $m=1,2,\cdots$, 令
\begin{align*}
\delta_m\triangleq\frac{1}{m+2},\qquad
\varepsilon_m\triangleq\frac{1}{2(m+2)},\qquad
V_m\triangleq(1-\delta_m,1+\delta_m).
\end{align*}
取 $f=\chi_{V_m}\in L^1(U)$, 令
\begin{align*}
U_{\varepsilon_m}
\triangleq (\varepsilon_m,2-\varepsilon_m),\quad
\varphi_m(x)
\triangleq 
\begin{cases}
\int_U
\eta_{\varepsilon_m}(x-y)\chi_{V_m}(y)\,\mathrm{d}y,
& x\in U_{\varepsilon_m},\\
0, & x\in\mathbb{R}\setminus U_{\varepsilon_m}.
\end{cases}
\end{align*}
由\hyperref[theorem:appC-光滑算子的性质]{光滑算子的性质\ref{theorem:appC-光滑算子的性质-1}\ref{theorem:appC-光滑算子的性质-6}}, 有
$\varphi_m\in C^\infty(U_{\varepsilon_m})$且
\begin{align}\label{eqFIJOUHJIOWjfiowjw38928ej234}
\operatorname{spt}\varphi_m
\subseteq
\overline{V_m}+[-\varepsilon_m,\varepsilon_m]
=[1-\delta_m-\varepsilon_m,1+\delta_m+\varepsilon_m]=\left[1-\frac{3}{2(m+2)},1+\frac{3}{2(m+2)}\right]\subseteq K\triangleq\left[\frac12,\frac32\right].
\end{align}
由于$K\subseteq U$,因而$\varphi_m\in C_c^\infty(U).$
由 $\eta_{\varepsilon_m}\geqslant0$ 和
$0\leqslant\chi_{V_m}\leqslant1$, 有
\begin{align*}
0\leqslant\varphi_m(x)
\leqslant\int_U\eta_{\varepsilon_m}(x-y)\,\mathrm{d}y
\leqslant\int_{\mathbb{R}}
\eta_{\varepsilon_m}(x-y)\,\mathrm{d}y
=1,\quad \forall x\in\mathbb{R}.
\end{align*}
又因为 $1\in U_{\varepsilon_m}$ 且
$\operatorname{spt}\eta_{\varepsilon_m}
\subseteq[-\varepsilon_m,\varepsilon_m]
\subset(-\delta_m,\delta_m)$,所以
\begin{align*}
\varphi_m(1)
=\int_{1-\delta_m}^{1+\delta_m}
\eta_{\varepsilon_m}(1-y)\,\mathrm{d}y
=\int_{-\delta_m}^{\delta_m}
\eta_{\varepsilon_m}(z)\,\mathrm{d}z
=\int_{\mathbb{R}}\eta_{\varepsilon_m}(z)\,\mathrm{d}z
=1.
\end{align*}

对于任意固定的 $x\neq1$, 当 $m$ 充分大时, 有$|x-1|>\frac{3}{2(m+2)}.$由\eqref{eqFIJOUHJIOWjfiowjw38928ej234}式可知, 此时 $\varphi_m(x)=0$. 因而
\begin{align}\label{eqJIOUJFIOIOW82u582382380234824890}
\lim\limits_{m\to\infty}\varphi_m(x)=0,
\qquad x\in\mathbb{R}\setminus\{1\}.
\end{align}
特别地, $\varphi_m\to0$ 在 $U$ 上几乎处处成立.

现在\eqref{align:ch05-弱导例2-点值恒等式}式中取
$\phi=\varphi_m$, 得
\begin{align}
1=\int_0^2v(x)\varphi_m(x)\,\mathrm{d}x
-\int_0^1\varphi_m(x)\,\mathrm{d}x.
\label{align:ch05-弱导例2-矛盾恒等式}
\end{align}
由于 $v\in L_{\mathrm{loc}}^1(U)$, 而 $K$ 是 $U$ 中的紧集, 有
\begin{align*}
|v(x)\varphi_m(x)|
\leqslant|v(x)|\chi_K(x),\qquad
\int_U|v(x)|\chi_K(x)\,\mathrm{d}x
=\int_K|v(x)|\,\mathrm{d}x<\infty.
\end{align*}
又由\eqref{eqJIOUJFIOIOW82u582382380234824890}式知$v\varphi_m\to0$在$U$上几乎处处, 由\hyperref[Real Analysis-theorem:控制收敛定理]{控制收敛定理}得
\begin{align*}
\lim\limits_{m\to\infty}
\int_0^2v(x)\varphi_m(x)\,\mathrm{d}x=0.
\end{align*}
同时,由$0\leqslant\varphi_m\leqslant1$ 及\eqref{align:ch05-弱导例2-点值恒等式}式得
\begin{align*}
0\leqslant\int_0^1\varphi_m(x)\,\mathrm{d}x
\leqslant\delta_m+\varepsilon_m
=\frac{3}{2(m+2)}
\longrightarrow0,\quad m\to \infty.
\end{align*}
因此, 在\eqref{align:ch05-弱导例2-矛盾恒等式}式中令
$m\to\infty$, 得$1=0$矛盾!所以不存在满足弱导数定义的
$v\in L_{\mathrm{loc}}^1(U)$, 即 $u$ 的弱导数不存在.
\end{solution}

\subsection{Sobolev空间的定义}\label{subsection:ch05-Sobolev空间定义}

\begin{definition}\label{definition:ch05-Wkp}
固定正整数 $k$ 与 $1\leqslant p\leqslant \infty$，设 $U \subseteq \mathbb{R}^n$ 是一个非空开集,\textbf{Sobolev空间}$W^{k,p}(U)$由所有局部可和函数 $u:U\to\mathbb{R}$ 组成，其中对每个满足 $|\alpha|\leqslant k$ 的多重指标 $\alpha$，$D^{\alpha}u$ 在弱意义下存在且属于 $L^{p}(U)$。
\end{definition}
\begin{remark}
类似\cref{definition:local_integrable_2026},此后我们将 $W^{k,p}(U)$ 中几乎处处相等的函数视为等同。
\end{remark}

\begin{definition}\label{definition:ch05-Wkp范数}
固定正整数 $k$ 与 $1\leqslant p\leqslant \infty$，设 $U \subseteq \mathbb{R}^n$ 是一个非空开集,若 $u\in W^{k,p}(U)$，定义其范数为
\begin{align*}
\|u\|_{W^{k,p}(U)}\triangleq
\begin{cases}
\left(\sum_{|\alpha|\leqslant k}\int_{U}|D^{\alpha}u|^{p}\,\mathrm{d}x\right)^{\frac{1}{p}}, & 1\leqslant p<\infty,\\
\sum_{|\alpha|\leqslant k}\operatorname*{ess\,sup}_{U}|D^{\alpha}u|, & p=\infty.
\end{cases}
\end{align*}
\begin{enumerate}[(1)]
\item 设 $\{u_{m}\}_{m=1}^{\infty}\subset W^{k,p}(U)$，$u\in W^{k,p}(U)$。我们称 $u_{m}$ 在 $W^{k,p}(U)$ 中收敛于 $u$，记作 $u_{m}\to u$ 于 $W^{k,p}(U)$，只要
\begin{align*}
\lim_{m\to\infty}\|u_{m}-u\|_{W^{k,p}(U)}=0.
\end{align*}
\item 我们写 $u_{m}\to u$ 于 $W^{k,p}_{\mathrm{loc}}(U)$，意指对每个 $V\subset \subset U$，都有 $u_{m}\to u$ 于 $W^{k,p}(V)$。
\end{enumerate}
我们记$W^{k,p}_{0}(U)$为 $C^{\infty}_{c}(U)$ 在 $W^{k,p}(U)$ 中的闭包。因此 $u\in W^{k,p}_{0}(U)$ 当且仅当存在函数 $u_{m}\in C^{\infty}_{c}(U)$ 使 $u_{m}\to u$ 于 $W^{k,p}(U)$。
当$p=2$时,我们通常记
\begin{align*}
H^{k}(U)=W^{k,2}(U),\quad H^{k}_{0}(U)=W^{k,2}_{0}(U).
\end{align*}
\end{definition}
\begin{remark}
我们将在\cref{proposition:W1p空间函数基本性质}中看到，若 $n=1$ 且 $U$ 为 $\mathbb{R}^{1}$ 中的开区间，则 $u\in W^{1,p}(U)$ 当且仅当 $u$ 几乎处处等于一个绝对连续函数，且其（几乎处处存在的）常义导数属于 $L^{p}(U)$。然而如此简单的刻画仅对 $n=1$ 可用。一般地，一个函数可以属于某个Sobolev空间，却仍然不连续或无界。
\end{remark}
\begin{remark}
Sobolev空间 \(W^{k,p}(U)\) 的范数如此定义的原因:
\begin{enumerate}[(1)]
\item \textbf{对所有 \(|\alpha|\leqslant k\) 的导数求和。} 这既是为了保证范数的正定性（避免常数函数与零函数混淆），也是因为Sobolev空间不依附于特定边界条件。与H\"older空间不同，最高阶导数的 \(L^p\) 可积性无法控制低阶导数的可积性。必须对直到 \(k\) 阶的所有弱导数进行统一控制，才能构造出一个完备的Banach空间。
\item \textbf{当 \(1\leqslant p<\infty\) 时开 \(\frac{1}{p}\) 次幂。} 这是为了满足范数公理中的齐次性，即将 \(p\) 次齐次转化为线性齐次 \(\|\lambda u\|=|\lambda|\|u\|\)，同时保证三角不等式成立，使得该空间成为赋范线性空间。

\item \textbf{当 \(p=\infty\) 时使用本性上确界。} 因为 \(W^{k,\infty}\) 中的元素是几乎处处相等的等价类，如果使用普通上确界，改变零测集上的函数值就会改变范数。使用 \(\operatorname*{ess\,sup}\) 可以忽略零测集的影响，使得范数在商空间上良定义。

\item \textbf{\(W_0^{k,p}(U)\) 定义为 \(C_c^\infty(U)\) 的闭包。} 偏微分方程需要刻画“在边界上为零”的边界条件，但边界是零测集，无法直接取逐点值。利用\cref{theorem:ch05-零迹函数}可知$W_0^{k,p}(U)$就是迹算子的核空间,因此我们把 $W^{k,p}_{0}(U)$ 理解为由那些 $u\in W^{k,p}(U)$ 中满足“$D^{\alpha}u=0$ 于 $\partial U$”（对一切 $|\alpha|\leqslant k-1$）的函数所组成。
\end{enumerate}
\end{remark}


\subsection{初等性质}\label{subsection:ch05-初等性质}

\begin{theorem}[弱导数的性质]\label{theorem:ch05-弱导数性质}
固定正整数 $k$ 与 $1\leqslant p\leqslant \infty$，设 $U \subseteq \mathbb{R}^n$ 是一个非空开集,$u,v\in W^{k,p}(U)$，$|\alpha|\leqslant k$。则
\begin{enumerate}[(1)]
\item\label{theorem:ch05-弱导数性质-1} $D^{\alpha}u\in W^{k-|\alpha|,p}(U)$，且对一切满足 $|\alpha|+|\beta|\leqslant k$ 的多重指标 $\alpha,\beta$，有 $D^{\beta}(D^{\alpha}u)=D^{\alpha}(D^{\beta}u)=D^{\alpha+\beta}u$。

\item\label{theorem:ch05-弱导数性质-2} 对每个 $\lambda,\mu\in\mathbb{R}$，$\lambda u+\mu v\in W^{k,p}(U)$，且 $D^{\alpha}(\lambda u+\mu v)=\lambda D^{\alpha}u+\mu D^{\alpha}v$，$|\alpha|\leqslant k$。

\item\label{theorem:ch05-弱导数性质-3} 若 $V$ 为 $U$ 的开子集，则 $u\in W^{k,p}(V)$。

\item\label{theorem:ch05-弱导数性质-4} 若 $\zeta\in C^{\infty}_{c}(U)$，则 $\zeta u\in W^{k,p}(U)$，且
\begin{align}
D^{\alpha}(\zeta u)=\sum_{\beta\leqslant\alpha}\binom{\alpha}{\beta}D^{\beta}\zeta\,D^{\alpha-\beta}u,
\label{align:ch05-leibniz}
\end{align}
此即 Leibniz 公式，其中 $\binom{\alpha}{\beta}=\dfrac{\alpha!}{\beta!(\alpha-\beta)!}$。
\end{enumerate}
\end{theorem}
\begin{proof}
\begin{enumerate}[(1)]
\item 先固定 $\phi\in C^{\infty}_{c}(U)$。则 $D^{\beta}\phi\in C^{\infty}_{c}(U)$，于是
\begin{align*}
\int_{U}D^{\alpha}u\,D^{\beta}\phi\,\mathrm{d}x
=(-1)^{|\alpha|}\int_{U}uD^{\alpha+\beta}\phi\,\mathrm{d}x=(-1)^{|\alpha|}(-1)^{|\alpha+\beta|}\int_{U}D^{\alpha+\beta}u\,\phi\,\mathrm{d}x=(-1)^{|\beta|}\int_{U}D^{\alpha+\beta}u\,\phi\,\mathrm{d}x.
\end{align*}
故在弱意义下 $D^{\beta}(D^{\alpha}u)=D^{\alpha+\beta}u$。

\item 任取测试函数 $\phi\in C^{\infty}_{c}(U)$。由弱导数的定义，有
\begin{align*}
\int_{U}uD^{\alpha}\phi\,\mathrm{d}x=(-1)^{|\alpha|}\int_{U}D^{\alpha}u\,\phi\,\mathrm{d}x,
\end{align*}
\begin{align*}
\int_{U}vD^{\alpha}\phi\,\mathrm{d}x=(-1)^{|\alpha|}\int_{U}D^{\alpha}v\,\phi\,\mathrm{d}x.
\end{align*}
利用积分的线性性质，计算 $\lambda u+\mu v$ 与 $D^{\alpha}\phi$ 的积分：
\begin{align*}
\int_{U}(\lambda u+\mu v)D^{\alpha}\phi\,\mathrm{d}x
&=\lambda\int_{U}uD^{\alpha}\phi\,\mathrm{d}x+\mu\int_{U}vD^{\alpha}\phi\,\mathrm{d}x\\
&=\lambda(-1)^{|\alpha|}\int_{U}D^{\alpha}u\,\phi\,\mathrm{d}x+\mu(-1)^{|\alpha|}\int_{U}D^{\alpha}v\,\phi\,\mathrm{d}x\\
&=(-1)^{|\alpha|}\int_{U}(\lambda D^{\alpha}u+\mu D^{\alpha}v)\phi\,\mathrm{d}x.
\end{align*}
由于 $\phi\in C^{\infty}_{c}(U)$ 是任意的，根据弱导数的定义，在弱意义下
\begin{align*}
D^{\alpha}(\lambda u+\mu v)=\lambda D^{\alpha}u+\mu D^{\alpha}v.
\end{align*}
此外，因为 $u,v\in W^{k,p}(U)$，所以 $\lambda u+\mu v\in L^{p}(U)$，且对一切 $|\alpha|\leqslant k$，$D^{\alpha}(\lambda u+\mu v)$ 作为 $L^{p}(U)$ 中元素的线性组合，依然属于 $L^{p}(U)$。因此 $\lambda u+\mu v\in W^{k,p}(U)$。

\item 因为 $u\in W^{k,p}(U)$，所以 $u\in L^{p}(U)$ 且对一切 $|\alpha|\leqslant k$，$D^{\alpha}u\in L^{p}(U)$。
由于 $V$ 为 $U$ 的开子集，限制在 $V$ 上的积分满足
\begin{align*}
\int_{V}|u|^{p}\,\mathrm{d}x\leqslant\int_{U}|u|^{p}\,\mathrm{d}x<\infty,
\end{align*}
同理 $\int_{V}|D^{\alpha}u|^{p}\,\mathrm{d}x<\infty$。因此 $u\in L^{p}(V)$ 且 $D^{\alpha}u\in L^{p}(V)$。

下面验证弱导数关系在 $V$ 上依然成立。任取函数 $\phi\in C^{\infty}_{c}(V)$。在 $U\setminus V$ 上令 $\phi\equiv 0$，将 $\phi$ 零延拓为 $U$ 上的函数，此时仍有 $\phi\in C^{\infty}_{c}(U)$。
因为 $u\in W^{k,p}(U)$，所以
\begin{align*}
\int_{U}uD^{\alpha}\phi\,\mathrm{d}x=(-1)^{|\alpha|}\int_{U}D^{\alpha}u\,\phi\,\mathrm{d}x.
\end{align*}
由于 $\phi$ 在 $U\setminus V$ 上恒为 $0$，上述积分在 $U\setminus V$ 上为 $0$，因此上式等价于
\begin{align*}
\int_{V}uD^{\alpha}\phi\,\mathrm{d}x=(-1)^{|\alpha|}\int_{V}D^{\alpha}u\,\phi\,\mathrm{d}x.
\end{align*}
这对任意 $\phi\in C^{\infty}_{c}(V)$ 成立，故在 $V$ 上 $D^{\alpha}u$ 依然是 $u$ 的弱导数。综上，$u\in W^{k,p}(V)$。

\item 我们对 $|\alpha|$ 作归纳来证式~\eqref{align:ch05-leibniz}。先设 $|\alpha|=1$。任取 $\phi\in C^{\infty}_{c}(U)$，则由可微函数的Leibniz公式知
\begin{align*}
D^{\alpha}\left( \zeta \phi \right) =\zeta D^{\alpha}\phi +\left( D^{\alpha}\zeta \right) \phi \Longrightarrow \zeta uD^{\alpha}\phi =uD^{\alpha}\left( \zeta \phi \right) -u\left( D^{\alpha}\zeta \right) \phi .
\end{align*}
从而
\begin{align*}
\int_{U}\zeta u\,D^{\alpha}\phi\,\mathrm{d}x
=\int_{U}uD^{\alpha}(\zeta\phi)-u(D^{\alpha}\zeta)\phi\,\mathrm{d}x=-\int_{U}(\zeta D^{\alpha}u+uD^{\alpha}\zeta)\phi\,\mathrm{d}x.
\end{align*}
故 $D^{\alpha}(\zeta u)=\zeta D^{\alpha}u+uD^{\alpha}\zeta$，即得所需。

现设 $l<k$，且\eqref{align:ch05-leibniz}式对一切 $|\alpha|\leqslant l$ 和一切函数 $\zeta$ 成立。取多重指标 $\alpha$ 满足 $|\alpha|=l+1$。则存在 $|\beta|=l$、$|\gamma|=1$ 使 $\alpha=\beta+\gamma$。设$\gamma$是第$j$个分量为1,其余分量为0的$n$维向量,则对如上的 $\phi$，又由\hyperref[Basis of Analytics-proposition:组合数基本性质]{组合数基本公式}可得
\begin{align*}
\binom{\beta}{\sigma -\gamma}+\binom{\beta}{\sigma}&=\left[ \binom{\beta _j}{\sigma _j-1}+\binom{\beta _j}{\sigma _j} \right] \prod_{i\ne j}{\binom{\beta _i}{\sigma _i}}=\binom{\beta _j+1}{\sigma _j}\prod_{i\ne j}{\binom{\beta _i}{\sigma _i}}
\\
&=\binom{\alpha _j}{\sigma _j}\prod_{i\ne j}{\binom{\alpha _i}{\sigma _i}}=\prod_{i=1}^n{\binom{\alpha _i}{\sigma _i}}=\binom{\alpha}{\sigma}.
\end{align*}
所以由归纳假设有
\begin{align*}
\int_{U}\zeta u\,D^{\alpha}\phi\,\mathrm{d}x
&=\int_U{\zeta u\,D^{\beta}\left( D^{\gamma}\phi \right) \,\mathrm{d}x}=\left( -1 \right) ^{\left| \beta \right|}\int_U{D^{\beta}\left( \zeta u \right) \,\left( D^{\gamma}\phi \right) \,\mathrm{d}x}\\
&=(-1)^{|\beta|}\int_{U}\sum_{\sigma\leqslant\beta}\binom{\beta}{\sigma}D^{\sigma}\zeta\,D^{\beta-\sigma}u\,D^{\gamma}\phi\,\mathrm{d}x
\\
&=(-1)^{|\beta|+|\gamma|}\int_{U}\sum_{\sigma\leqslant\beta}\binom{\beta}{\sigma}D^{\gamma}(D^{\sigma}\zeta\,D^{\beta-\sigma}u)\phi\,\mathrm{d}x
\\
&=\left( -1 \right) ^{|\alpha |}\int_U{\sum_{\sigma \leqslant \beta}{\binom{\beta}{\sigma}\left[ D^{\sigma +\gamma}\zeta \,D^{\alpha -\left( \sigma +\gamma \right)}u+D^{\sigma}\zeta \,D^{\alpha -\sigma}u \right] \phi \,\mathrm{d}x}}\quad\ (|\gamma|=1)
\\
&=\left( -1 \right) ^{|\alpha |}\int_U{\left[ \sum_{\rho -\gamma \leqslant \beta}{\binom{\beta}{\rho -\gamma}D^{\rho}\zeta \,D^{\alpha -\rho}u}+\sum_{\sigma \leqslant \beta}{\binom{\beta}{\sigma}D^{\sigma}\zeta \,D^{\alpha -\sigma}u} \right] \phi \,\mathrm{d}x}\quad\text{（其中}\ \rho=\sigma+\gamma\text{）}
\\
&=\left( -1 \right) ^{|\alpha |}\int_U{\left[ \sum_{\rho \leqslant \alpha}{\binom{\beta}{\rho -\gamma}D^{\rho}\zeta \,D^{\alpha -\rho}u}+\sum_{\sigma \leqslant \alpha}{\binom{\beta}{\sigma}D^{\sigma}\zeta \,D^{\alpha -\sigma}u} \right] \phi \,\mathrm{d}x}
\\
&=\left( -1 \right) ^{|\alpha |}\int_U{\left[ \sum_{\sigma \leqslant \alpha}{\binom{\beta}{\sigma -\gamma}D^{\sigma}\zeta \,D^{\alpha -\sigma}u}+\sum_{\sigma \leqslant \alpha}{\binom{\beta}{\sigma}D^{\sigma}\zeta \,D^{\alpha -\sigma}u} \right] \phi \,\mathrm{d}x}
\\
&=\left( -1 \right) ^{|\alpha |}\int_U{\left[ \sum_{\sigma \leqslant \alpha}{\binom{\alpha}{\sigma}D^{\sigma}\zeta \,D^{\alpha -\sigma}u} \right] \phi \,\mathrm{d}x}.
\end{align*}
\end{enumerate}
\end{proof}

\begin{example}\label{example:ch05-径向奇异}
取 $U=B^{0}(0,1)$，即 $\mathbb{R}^{n}$ 中的开单位球，令
\begin{align*}
u(x)=|x|^{-\alpha}\qquad (x\in U,\ x\neq 0).
\end{align*}
对哪些 $\alpha>0$、$n$、$p$ 值，$u$ 属于 $W^{1,p}(U)$？
\end{example}
\begin{solution}

我们断言:对于 $\alpha>0$, 整数 $n\geqslant1$ 以及
$1\leqslant p\leqslant\infty$, 有
\begin{align*}
u\in W^{1,p}(U)
\Longleftrightarrow
1\leqslant p<n
\ \text{且}\
0<\alpha<\frac{n-p}{p}.
\end{align*}
特别地, 对每个 $p\geqslant n$, 都有 $u\notin W^{1,p}(U)$.
当 $n=1$ 时, 对任何 $\alpha>0$ 和 $1\leqslant p\leqslant\infty$, 均有
$u\notin W^{1,p}(U)$.

由于单点集的 Lebesgue 测度为零, 可以令 $u(0)=0$,而不影响结论. 
记 $\alpha(n)$ 为 $\mathbb{R}^n$ 中开单位球的体积.首先, $u$ 在 $U\setminus\{0\}$ 上光滑, 其在通常意义下的梯度为
\begin{align*}
g(x)\triangleq-\alpha\frac{x}{|x|^{\alpha+2}},
\quad x\neq0.
\end{align*}
令 $g(0)=0$, 则
\begin{align*}
g_i(x)=-\alpha\frac{x_i}{|x|^{\alpha+2}},
\quad
|g(x)|=\frac{\alpha}{|x|^{\alpha+1}},
\quad x\neq0.
\end{align*}

先考虑 $1\leqslant p<\infty$. 由\hyperref[theorem:appC-极坐标]{极坐标公式}, 有
\begin{align*}
\int_U{|u(x)|^p\,\mathrm{d}x}&=\int_0^1{\left( \int_{\partial B\left( 0,r \right)}{\frac{1}{\left| x \right|^{\alpha p}}\,\mathrm{d}S} \right) \mathrm{d}r}=\int_0^1{\left( \int_{\partial B\left( 0,r \right)}{\frac{1}{r^{\alpha p}}\,\mathrm{d}S} \right) \mathrm{d}r}
\\
&=\int_0^1{n\alpha \left( n \right) r^{n-1}\cdot \frac{1}{r^{\alpha p}}\mathrm{d}r}=n\alpha (n)\int_0^1{r^{n-1-\alpha p}\,\mathrm{d}r},
\end{align*}
\begin{align*}
\int_U{\left| g\left( x \right) \right|^p\,\mathrm{d}x}&=\int_0^1{\left( \int_{\partial B\left( 0,r \right)}{\frac{\left| \alpha \right|^p}{\left| x \right|^{\left( \alpha +1 \right) p}}\,\mathrm{d}S} \right) \mathrm{d}r}=\int_0^1{\left( \int_{\partial B\left( 0,r \right)}{\frac{\left| \alpha \right|^p}{r^{\left( \alpha +1 \right) p}}\,\mathrm{d}S} \right) \mathrm{d}r}
\\
&=\int_0^1{n\alpha \left( n \right) r^{n-1}\cdot \frac{\left| \alpha \right|^p}{r^{\left( \alpha +1 \right) p}}\mathrm{d}r}=n\alpha \left( n \right) \left| \alpha \right|^p\int_0^1{r^{\left( n-1 \right) -\left( \alpha +1 \right) p}\mathrm{d}r}.
\end{align*}
由于
\begin{align*}
\int_0^1r^\beta\,\mathrm{d}r<\infty
\Longleftrightarrow
\beta>-1,
\end{align*}
并且 $\alpha>0$, 所以
\begin{align}
u\in L^p(U)
\Longleftrightarrow
\alpha p<n,
\quad
g\in L^p(U)
\Longleftrightarrow
(\alpha+1)p<n.
\label{align:ch05-径向奇异-梯度可积性}
\end{align}

{\heiti 充分性:}
假设$1\leqslant p<n
\ \text{且}\
0<\alpha<\frac{n-p}{p}$,则$(\alpha+1)p<n$. 由于 $p\geqslant1$, 有
\begin{align*}
\alpha+1\leqslant(\alpha+1)p<n.
\end{align*}
因此 $\alpha<n-1$, 特别地, $\alpha<n$. 再次由\hyperref[theorem:appC-极坐标]{极坐标公式},得
\begin{align*}
&\int_U{|u(x)|\,\mathrm{d}x}=\int_0^1{\left( \int_{\partial B\left( 0,r \right)}{\frac{1}{\left| x \right|^{\alpha}}\,\mathrm{d}S} \right) \mathrm{d}r}=\int_0^1{\left( \int_{\partial B\left( 0,r \right)}{\frac{1}{r^{\alpha}}\,\mathrm{d}S} \right) \mathrm{d}r}
\\
&=\int_0^1{n\alpha \left( n \right) r^{n-1}\cdot \frac{1}{r^{\alpha}}\mathrm{d}r}=n\alpha (n)\int_0^1{r^{n-1-\alpha}\,\mathrm{d}r}=\frac{n\alpha \left( n \right)}{n-\alpha}<\infty ,
\end{align*}
\begin{align*}
&\int_U{\left| g\left( x \right) \right|\,\mathrm{d}x}=\int_0^1{\left( \int_{\partial B\left( 0,r \right)}{\frac{\alpha}{\left| x \right|^{\alpha +1}}\,\mathrm{d}S} \right) \mathrm{d}r}=\int_0^1{\left( \int_{\partial B\left( 0,r \right)}{\frac{\alpha}{r^{\alpha +1}}\,\mathrm{d}S} \right) \mathrm{d}r}
\\
&=\int_0^1{n\alpha \left( n \right) r^{n-1}\cdot \frac{\alpha}{r^{\alpha +1}}\mathrm{d}r}=n\alpha \left( n \right) \alpha \int_0^1{r^{n-2-\alpha}\mathrm{d}r}=\frac{n\alpha \left( n \right) \alpha}{n-1-\alpha}<\infty .
\end{align*}
故 $u\in L^1(U)$ 且 $g\in L^1(U)$.

任取 $\phi\in C_c^\infty(U)$, 固定 $i\in\{1,\cdots,n\}$. 因为 $\phi$ 及其一阶导数连续且具有紧支集, 所以它们在 $U$ 上有界. 从而
\begin{gather*}
\int_U|uD_i\phi|\,\mathrm{d}x
\leqslant
\|D_i\phi\|_{L^\infty(U)}
\int_U|u|\,\mathrm{d}x<\infty,\\
\int_U|\phi g_i|\,\mathrm{d}x
\leqslant
\|\phi\|_{L^\infty(U)}
\int_U|g_i|\,\mathrm{d}x\leqslant
\|\phi\|_{L^\infty(U)}
\int_U|g|\,\mathrm{d}x<\infty.
\end{gather*}
这说明 $uD_i\phi,\phi g_i\in L^1(U)$.

对 $0<\varepsilon<1$, 令
\begin{align*}
\Omega_\varepsilon
\triangleq\{x\in\mathbb{R}^n:\varepsilon<|x|<1\}.
\end{align*}
由于$\phi\in C_c^\infty(U)$, 可将其光滑地延伸到 $\partial U$ 并取值为零. 在 $\Omega_\varepsilon$ 上应用
\hyperref[Basis of Analytics-theorem:Green恒等式-3]{Green恒等式\ref{Basis of Analytics-theorem:Green恒等式-3}}, 得
\begin{align*}
\int_{\Omega_\varepsilon}uD_i\phi\,\mathrm{d}x
&=-\int_{\Omega_\varepsilon}\phi g_i\,\mathrm{d}x
+\int_{\partial B^{0}(0,\varepsilon)}u\phi\nu_i\,\mathrm{d}S
+\int_{\partial U}u\phi\mu_i\,\mathrm{d}S.
\end{align*}
其中, $\nu=(\nu_1,\cdots,\nu_n)$ 是 $\Omega_\varepsilon$ 在内边界上的单位外法向量, 也就是小球$B^{0}(0,\varepsilon)$的单位内法向量, 而 $\mu=(\mu_1,\cdots,\mu_n)$ 是 $U$ 的单位外法向量. 具体地,有
\begin{align*}
\nu (x)=-\frac{x}{\varepsilon}\,\,\left( |x|=\varepsilon \right) ,\quad \mu (x)=x\,\,\left( |x|=1 \right) .
\end{align*}
由$\phi\in C_c^\infty(U)$知$\phi$ 在 $\partial U$ 上为零, 外边界积分消失, 因而
\begin{align}\label{eq::JIOJIOWJ22221}
\int_{\Omega_\varepsilon}uD_i\phi\,\mathrm{d}x
=-\int_{\Omega_\varepsilon}\phi g_i\,\mathrm{d}x
+\int_{\partial B^{0}(0,\varepsilon)}u\phi\nu_i\,\mathrm{d}S.
\end{align}
由 $uD_i\phi,\phi g_i\in L^1(U)$ 以及
$|B^{0}(0,\varepsilon)|=\alpha(n)\varepsilon^n\to0\,(\varepsilon \to 0^+)$, 根据\hyperref[Real Analysis-theorem:积分的绝对连续性]{积分的绝对连续性},当$\varepsilon \to 0^+$时,有
\begin{align*}
\left|
\int_{B^{0}(0,\varepsilon)}uD_i\phi\,\mathrm{d}x
\right|
&\leqslant
\int_{B^{0}(0,\varepsilon)}|uD_i\phi|\,\mathrm{d}x
\longrightarrow0,\\
\left|
\int_{B^{0}(0,\varepsilon)}\phi g_i\,\mathrm{d}x
\right|
&\leqslant
\int_{B^{0}(0,\varepsilon)}|\phi g_i|\,\mathrm{d}x
\longrightarrow0.
\end{align*}
又由于球面$\partial B^{0}(0,\varepsilon)$是零Lebesgue测度集,因此
\begin{gather*}
\lim_{\varepsilon \rightarrow 0^+} \int_{\Omega _{\varepsilon}}{uD_i\phi \,\mathrm{d}x}=\int_U{uD_i\phi \,\mathrm{d}x}-\lim_{\varepsilon \rightarrow 0^+} \int_{B^0(0,\varepsilon )}{uD_i\phi \,\mathrm{d}x}=\int_U{uD_i\phi \,\mathrm{d}x},\quad 
\\
\lim_{\varepsilon \rightarrow 0^+} \int_{\Omega _{\varepsilon}}{\phi g_i\,\mathrm{d}x}=\int_U{\phi g_i\,\mathrm{d}x}-\lim_{\varepsilon \rightarrow 0^+} \int_{B^0(0,\varepsilon )}{\phi g_i\,\mathrm{d}x}=\int_U{\phi g_i\,\mathrm{d}x}.
\end{gather*}
另一方面, 利用 $|\nu_i|\leqslant1$ 以及 $\alpha<n-1$, 得
\begin{align*}
\left|
\int_{\partial B^{0}(0,\varepsilon)}u\phi\nu_i\,\mathrm{d}S
\right|
\leqslant
\|\phi\|_{L^\infty(U)}
\int_{\partial B^{0}(0,\varepsilon)}
\varepsilon^{-\alpha}\,\mathrm{d}S
=n\alpha(n)\|\phi\|_{L^\infty(U)}
\varepsilon^{n-1-\alpha}
\longrightarrow0.
\end{align*}
于是, 在 \eqref{eq::JIOJIOWJ22221} 中令 $\varepsilon\to0^+$, 可得
\begin{align*}
\int_UuD_i\phi\,\mathrm{d}x
=-\int_Ug_i\phi\,\mathrm{d}x,
\quad \forall \phi\in C_c^\infty(U).
\end{align*}
由于 $i$ 任意, 故$g$ 就是 $u$ 在整个 $U$ 上的一阶弱导数.
又由充分性假设知
\begin{align*}
\alpha p<(\alpha+1)p<n,
\end{align*}
再由\eqref{align:ch05-径向奇异-梯度可积性}式, 有
$u\in L^p(U)$ 且 $g\in L^p(U)$. 因此
$u\in W^{1,p}(U)$.

{\heiti 必要性:}
假设$1\leqslant p<\infty,u\in W^{1,p}(U)$, 记其一阶弱导数为
$v=(v_1,\cdots,v_n)\in L^p(U)$. 对任意
$\psi\in C_c^\infty(U\setminus\{0\})$, 将其在原点处定义为零后, 它仍是
$C_c^\infty(U)$ 中的测试函数. 由弱导数的定义知
\begin{align*}
\int_UuD_i\psi\,\mathrm{d}x
=-\int_Uv_i\psi\,\mathrm{d}x.
\end{align*}
另一方面, 由于 $\psi$ 的支集避开原点, 而 $u$ 在
$U\setminus\{0\}$ 上光滑,由分部积分给出
\begin{align*}
\int_{U\setminus \{0\}} uD_i\psi\,\mathrm{d}x
=-\int_{U\setminus \{0\}} g_i\psi\,\mathrm{d}x.
\end{align*}
因此
\begin{align*}
\int_{U\setminus\{0\}}(v_i-g_i)\psi\,\mathrm{d}x=0,
\qquad
\psi\in C_c^\infty(U\setminus\{0\}).
\end{align*}
由\cref{proposition:整体可积与局部可积的关系}知$v_i\in L^1(U)\subseteq L^1_{\mathrm{loc}}(U\setminus\{0\})$,又由$g_i\in C(U\setminus \{0\})$知$g_i\in L^1_{\mathrm{loc}}(U\setminus\{0\})$,故根据\hyperref[lemma:变分法基本原理]{变分法基本原理},有
\begin{align*}
v_i=g_i,
\quad \text{在 }U\setminus\{0\}\text{ 上几乎处处成立}.
\end{align*}
单点集 $\{0\}$ 的测度为零, 所以 $v=g$ 在整个 $U$ 上几乎处处成立. 从而
\begin{align*}
\int_U|g(x)|^p\,\mathrm{d}x
=\int_U|v(x)|^p\,\mathrm{d}x<\infty.
\end{align*}
由\eqref{align:ch05-径向奇异-梯度可积性}式, 必须有
\begin{align*}
(\alpha+1)p<n\Longrightarrow  1\leqslant p<n
\ \text{且}\
0<\alpha<\frac{n-p}{p}.
\end{align*}

下面考虑$p=\infty$的情形.设 $u\in W^{1,\infty}(U)$, 记其一阶弱偏导数为 $v_i$, $i=1,\cdots,n$. 则
\begin{align*}
u\in L^\infty(U),\qquad
v_i\in L^\infty(U),
\qquad i=1,\cdots,n.
\end{align*}
由于 $|U|=\alpha(n)<\infty$, 有
\begin{align*}
\int_U|u(x)|^n\,\mathrm{d}x
&\leqslant |U|\|u\|_{L^\infty(U)}^n<\infty,\\
\int_U|v_i(x)|^n\,\mathrm{d}x
&\leqslant |U|\|v_i\|_{L^\infty(U)}^n<\infty,
\qquad i=1,\cdots,n.
\end{align*}
因此 $u$ 及其所有一阶弱偏导数均属于 $L^n(U)$, 从而$u\in W^{1,n}(U).$
由于 $n$ 是有限的正整数, 可以应用前面已经证明的$p$有限的情形, 得$(\alpha+1)n<n.$但 $\alpha>0$ 且 $n\geqslant1$, 所以 $(\alpha+1)n>n$, 矛盾! 
\end{solution}

\begin{example}\label{example:ch05-稠密点奇异}
设 $\{r_{k}\}_{k=1}^{\infty}$ 为 $U=B^{0}(0,1)$ 的一个可数稠密子集。写
\begin{align*}
u(x)=\sum_{k=1}^{\infty}\frac{1}{2^{k}}|x-r_{k}|^{-\alpha}\qquad (x\in U).
\end{align*}
则当 $0<\alpha<\dfrac{n-p}{p}$ 时,有$u\in W^{1,p}(U)$，但 $u$ 在 $U$ 的每个开子集上都无界。
\end{example}
\begin{remark}
这个例子说明了一个基本事实：尽管属于Sobolev空间的函数 $u$ 具有某些光滑性态，它在其他方面仍可能性态相当恶劣。
\end{remark}
\begin{proof}
由条件 $0<\alpha<\frac{n-p}{p}$, 有
\begin{align*}
\alpha p<(\alpha+1)p<n.
\end{align*}
记 $\alpha(n)$ 为 $\mathbb{R}^n$ 中开单位球的体积. 对每个 $k\in\mathbb{N}$, 定义
\begin{align*}
f_k(x)\triangleq
\begin{cases}
|x-r_k|^{-\alpha}, & x\in U\setminus\{r_k\},\\
0, & x=r_k,
\end{cases}
\end{align*}
并令
\begin{align*}
u_m(x)\triangleq\sum\limits_{k=1}^{m}\frac{1}{2^k}f_k(x),
\qquad m\in\mathbb{N}.
\end{align*}

先证明级数几乎处处收敛到有限值, 且其和属于 $L^p(U)$.
由于 $x,r_k\in B^{0}(0,1)$, 有
\begin{align*}
|x-r_k|\leqslant|x|+|r_k|<2.
\end{align*}
作平移变换 $y=x-r_k$, 利用 $U-r_k\subseteq B^{0}(0,2)$, 得
\begin{align*}
\|f_k\|_{L^p(U)}^p
=\int_{U-r_k}|y|^{-\alpha p}\,\mathrm{d}y\leqslant\int_{B^{0}(0,2)}|y|^{-\alpha p}\,\mathrm{d}y.
\end{align*}
由\hyperref[theorem:appC-极坐标]{极坐标公式},得
\begin{align}\label{eq::JIFJIJOIWjjjjiij}
\int_{B^0(0,2)}{|y|^{-\alpha p}\,\mathrm{d}y}=\int_0^2{\int_{\partial B(0,1)}{r^{-\alpha p}\mathrm{d}S\,\mathrm{d}r}}=n\alpha (n)\int_0^2{r^{n-1-\alpha p}\,\mathrm{d}r}=\frac{n\alpha (n)2^{n-\alpha p}}{n-\alpha p}<\infty .
\end{align}
因此,存在与 $k$ 无关的常数
$C_0\triangleq
\left(\frac{n\alpha(n)2^{n-\alpha p}}{n-\alpha p}\right)^{\frac{1}{p}},$
使得
\begin{align*}
\|f_k\|_{L^p(U)}\leqslant C_0.
\end{align*}
应用\hyperref[Real Analysis-theorem:实变函数--Minkowski不等式]{Minkowski不等式}, 得
\begin{align*}
\|u_m\|_{L^p(U)}
\leqslant
\sum\limits_{k=1}^{m}\frac{1}{2^k}\|f_k\|_{L^p(U)}
\leqslant C_0\sum\limits_{k=1}^{m}\frac{1}{2^k}
\leqslant C_0.
\end{align*}

因为 $f_k\geqslant0$, 所以 $\{u_m\}$ 逐点关于$m$单调递增,由单调有界定理,可令
\begin{align*}
w(x)\triangleq\lim\limits_{m\to\infty}u_m(x)\in[0,+\infty].
\end{align*}
由\hyperref[Real Analysis-theorem:Beppo Levi非负渐升列积分定理]{Levi定理},得
\begin{align*}
\int_Uw^p(x)\,\mathrm{d}x
=\lim\limits_{m\to\infty}\int_Uu_m^p(x)\,\mathrm{d}x
\leqslant C_0^p<\infty.
\end{align*}
因此 $w$ 几乎处处有限. 记
\begin{align}\label{eq:::jiijijiojiojwjiowfjiojfwio}
N\triangleq
\{r_k:k\in\mathbb{N}\}\cup\{x\in U:w(x)=+\infty\},
\end{align}
则 $N$ 为零测集. 选取代表元
\begin{align*}
u(x)\triangleq
\begin{cases}
w(x), & x\in U\setminus N,\\
0, & x\in N.
\end{cases}
\end{align*}
于是 $u\in L^p(U)$, 且 $u$ 在 $U\setminus N$ 上等于题设级数,即$u$等价于题设级数.

此外, 在 $U\setminus N$ 上,有
\begin{align*}
0\leqslant u_m\leqslant u,\qquad
u_m\longrightarrow u,\qquad
|u-u_m|^p\leqslant u^p.
\end{align*}
由$u\in L^p(U)$知$u^p\in L^1(U)$, 由\hyperref[Real Analysis-theorem:控制收敛定理]{控制收敛定理},得
\begin{align*}
\lim\limits_{m\to\infty}\|u_m-u\|_{L^p(U)}^p
=\lim\limits_{m\to\infty}
\int_U|u_m-u|^p\,\mathrm{d}x
=0.
\end{align*}

下面证明每个 $f_k$ 的弱梯度为
\begin{align*}
g_k(x)\triangleq
\begin{cases}
-\alpha\displaystyle\frac{x-r_k}{|x-r_k|^{\alpha+2}},
& x\in U\setminus\{r_k\},\\
0, & x=r_k.
\end{cases}
\end{align*}
由\cref{example:ch05-径向奇异}, 函数
\begin{align*}
h(y)=|y|^{-\alpha},
\qquad y\in B^{0}(0,1)\setminus\{0\},
\end{align*}
在原点处任意赋予有限值后属于 $W^{1,p}(B^{0}(0,1))$, 且其弱偏导数满足
\begin{align*}
D_i h(y)=-\alpha\frac{y_i}{|y|^{\alpha+2}},
\qquad y\neq0,\quad i=1,\cdots,n.
\end{align*}
任取 $\phi\in C_c^\infty(U)$, 将其在 $U$ 外定义为零, 得到 $\mathbb{R}^n$ 上的光滑函数. 因为 $U\subseteq B^{0}(r_k,2)$, 所以
\begin{align*}
\psi_k(y)\triangleq\phi(r_k+2y)\in C_c^\infty(B^{0}(0,1)).
\end{align*}
再令
\begin{align*}
\widetilde{\phi }\left( x \right) \triangleq \begin{cases}
\phi \left( x \right) ,&x\in U,\\
0,&x\in \mathbb{R} \backslash U.\\
\end{cases}
\end{align*}
则$\overline{\phi}\in C_c^\infty(U)\subseteq C_c^\infty(B^0(r_k,2))$.作变量替换 $x=r_k+2y,y=\frac{x-r_k}{2}$, 并利用链式法则,得
\begin{align*}
\int_U{f_k\left( x \right) D_{x_i}\phi \left( x \right) \,\mathrm{d}x}&=\int_{B^0(r_k,2)}{\left| x-r_k \right|^{-\alpha}D_{x_i}\widetilde{\phi }\left( x \right) \,\mathrm{d}x}\xlongequal{x=r_k+2y}\int_{B^0(0,1)}{2^{-\alpha}\left| y \right|^{-\alpha}D_{y_i}\psi _k\left( y \right) \cdot \frac{\mathrm{d}y_i}{\mathrm{d}x_i}\cdot \,2^n\,\mathrm{d}y}
\\
&=2^{n-\alpha}\int_{B^0(0,1)}{h\left( y \right) \cdot \frac{1}{2}D_{y_i}\psi _k\left( y \right) \,\mathrm{d}y}=2^{n-\alpha -1}\int_{B^0(0,1)}{h\left( y \right) D_{y_i}\psi _k\left( y \right) \,\mathrm{d}y}
\\
&=-2^{n-\alpha -1}\int_{B^0(0,1)}{\psi _k\left( y \right) D_ih\left( y \right) \,\mathrm{d}y}\xlongequal{y=\frac{x-r_k}{2}}-2^{-\alpha -1}\int_{B^0(r_k,2)}{\widetilde{\phi }_k\left( y \right) D_ih\left( \frac{x-r_k}{2} \right) \,\mathrm{d}x}
\\
&=-2^{-\alpha -1}\int_U{\phi _k\left( y \right) D_ih\left( \frac{x-r_k}{2} \right) \,\mathrm{d}x}=-2^{-\alpha -1}\int_U{\phi _k\left( y \right) \frac{-\alpha \cdot \frac{x_i-r_{k,i}}{2}}{\left| \frac{x-r_k}{2} \right|^{\alpha +2}}\,\mathrm{d}x}
\\
&=2^{-\alpha -1}\int_U{\phi _k\left( y \right) \cdot \left( -\alpha 2^{\alpha +1}\frac{x_i-r_{k,i}}{\left| x-r_k \right|^{\alpha +2}} \right) \,\mathrm{d}x}=\int_U{\phi _k\left( y \right) g_{k,i}\left( x \right) \,\mathrm{d}x}.
\end{align*}
因此 $g_{k,i}$ 确实是 $f_k$ 在 $U$ 上的弱偏导数.

又因为
\begin{align*}
|g_k(x)|=\alpha|x-r_k|^{-\alpha-1},
\qquad x\neq r_k,
\end{align*}
同样做平移变换$y=x-r_k$,利用 $U-r_k\subseteq B^{0}(0,2)$和\hyperref[theorem:appC-极坐标]{极坐标公式},得
\begin{align*}
\left\| g_k \right\| _{L^p\left( U \right)}^{p}&=\alpha ^p\int_U{\left| x-r_k \right|^{-(\alpha +1)p}\,\mathrm{d}x}=\alpha ^p\int_{U-r_k}{\left| y \right|^{-(\alpha +1)p}\,\mathrm{d}y}
\\
&\leqslant \alpha ^p\int_{B^0(0,2)}{\left| y \right|^{-(\alpha +1)p}\,\mathrm{d}y}=n\alpha (n)\alpha ^p\int_0^2{r^{n-1-(\alpha +1)p}\,\mathrm{d}r}
\\
&=\frac{n\alpha (n)\alpha ^p2^{n-(\alpha +1)p}}{n-(\alpha +1)p}<\infty .
\end{align*}
所以 $f_k\in W^{1,p}(U)$, 且存在与 $k$ 无关的常数$C_1\triangleq
\left(
\frac{n\alpha(n)\alpha^p2^{n-(\alpha+1)p}}
{n-(\alpha+1)p}
\right)^{\frac{1}{p}},$
使得
\begin{align*}
\|g_k\|_{L^p(U)}\leqslant C_1.
\end{align*}
再令
\begin{align*}
v_m\triangleq
\sum\limits_{k=1}^{m}\frac{1}{2^k}g_k,
\end{align*}
则由\hyperref[theorem:ch05-弱导数性质]{弱导数的性质\ref{theorem:ch05-弱导数性质-2}}, $u_m\in W^{1,p}(U)$, 且 $Du_m=v_m$.
当 $\ell>m$ 时, 再由\hyperref[Real Analysis-theorem:实变函数--Minkowski不等式]{Minkowski不等式},得
\begin{align*}
\|v_\ell-v_m\|_{L^p(U)}
\leqslant
\sum\limits_{k=m+1}^{\ell}
\frac{1}{2^k}\|g_k\|_{L^p(U)}
\leqslant C_1\sum\limits_{k=m+1}^{\ell}\frac{1}{2^k}
\leqslant\frac{C_1}{2^m}.
\end{align*}
因此 $\{v_m\}$ 是 $L^p(U)$ 中的 Cauchy 列. 由$L^p(U)$空间的完备性, 存在 $v\in L^p(U)$, 使得
\begin{align*}
v_m\longrightarrow v
\qquad\text{于 }L^p(U).
\end{align*}

下面验证 $v$ 就是 $u$ 的弱导数. 任取 $\phi\in C_c^\infty(U)$, 固定 $i\in\{1,\cdots,n\}$. 对每个 $m$, 有
\begin{align}
\int_Uu_mD_i\phi\,\mathrm{d}x
=-\int_Uv_{m,i}\phi\,\mathrm{d}x.
\label{align:ch05-稠密点奇异-部分和弱导数}
\end{align}
记 $q$ 为 $p$ 的共轭指数, 其中 $p=1$ 时取 $q=\infty$.
因为 $\phi$ 及其导数光滑且具有紧支集, 所以
$\phi,D_i\phi\in L^q(U)$. 由\hyperref[Real Analysis-theorem:实变函数--Hölder不等式]{Hölder不等式},得
\begin{align*}
\left|
\int_U(u_m-u)D_i\phi\,\mathrm{d}x
\right|
&\leqslant
\|u_m-u\|_{L^p(U)}\|D_i\phi\|_{L^q(U)}
\longrightarrow0,\\
\left|
\int_U(v_{m,i}-v_i)\phi\,\mathrm{d}x
\right|
&\leqslant
\|v_m-v\|_{L^p(U)}
\|\phi\|_{L^q(U)}
\longrightarrow0.
\end{align*}
因此, 在\eqref{align:ch05-稠密点奇异-部分和弱导数}中令 $m\to\infty$, 得
\begin{align*}
\int_UuD_i\phi\,\mathrm{d}x
=-\int_Uv_i\phi\,\mathrm{d}x,
\qquad \phi\in C_c^\infty(U).
\end{align*}
所以 $Du=v\in L^p(U)$. 结合 $u\in L^p(U)$, 得
\begin{align*}
u\in W^{1,p}(U).
\end{align*}

最后, 证明 $u$ 在每个非空开子集上无界.
任取非空开集 $V\subseteq U$. 由 $\{r_k\}_{k=1}^{\infty}$ 的稠密性, 存在 $m$ 使得 $r_m\in V$.
选取 $\rho>0$, 使得 $B^{0}(r_m,\rho)\subseteq V$.

回顾\eqref{eq:::jiijijiojiojwjiowfjiojfwio}式中$N$的定义,它包含所有 $r_k$, 且在 $U\setminus N$ 上, $u$ 等于给定级数的和.
对于任意 $x\in V\setminus N$,有
\begin{align}\label{qeqjijiowjiojiowiwiwiwj22}
u(x)
=\sum\limits_{k=1}^{\infty}\frac{1}{2^k}|x-r_k|^{-\alpha}
\geqslant\frac{1}{2^m}|x-r_m|^{-\alpha}.
\end{align}
由于 $|N|=0$, 对任意 $0<t<\rho$, 均有
\begin{align*}
|B^{0}(r_m,t)\setminus N|
=\alpha(n)t^n>0\Longrightarrow B^{0}(r_m,t)\setminus N\neq \varnothing.
\end{align*}
因此可取$x_n\in B^0\left( r_m,\frac{1}{n\rho} \right) \setminus N\,(n=1,2,\cdots)$,则$x_n\to r_m\,(n\to \infty)$.又因为$\alpha>0$,所以\eqref{qeqjijiowjiojiowiwiwiwj22}式给出
\begin{align*}
\underset{n\rightarrow \infty}{\lim}u\left( x_n \right) \geqslant \frac{1}{2^m}\left| x_n-r_m \right|^{-\alpha}=+\infty .
\end{align*}
从而$u$ 在 $V$ 上无界. 由于 $V$ 任意, 故$u$ 在 $U$ 的每个非空开子集上都无界.
\end{proof}

\begin{proposition}\label{proposition:ch05--examoi-11}
设 $U\subseteq\mathbb{R}^n$ 为非空连通开集, $1\leqslant p\leqslant\infty$, 且 $u\in W^{1,p}(U)$ 满足
\begin{align*}
Du(x)=0,
\quad \text{a.e.}\,x\in U.
\end{align*}
则存在常数 $c\in\mathbb{R}$, 使得 $u(x)=c\,\,\text{a.e.}\,x\in U.$
\end{proposition}
\begin{proof}

由于 $u\in W^{1,p}(U)$,结合\cref{proposition:整体可积与局部可积的关系}有 $u\in L^p(U)
\subseteq L^p_{\mathrm{loc}}(U) \subseteq L^1_{\mathrm{loc}}(U)$.
又因为 $Du=0$ 几乎处处成立, 所以根据弱偏导数的定义, 对任意
$i\in\{1,\cdots,n\}$ 和 $\varphi\in C_c^\infty(U)$, 有
\begin{align}
\int_Uu(y)D_i\varphi(y)\,\mathrm{d}y
=-\int_UD_i u(y)\varphi(y)\,\mathrm{d}y
=0.
\label{eqwr:Njkojioejioio2}
\end{align}

$\mathbf{Step}\mathbf{1.}$先证明 $u$ 在每个闭包包含于 $U$ 的开球上几乎处处为常数.
任取开球 $B=B^{0}(a,r)$, 满足 $\overline{B}\subseteq U$.
由于 $\overline{B}$ 紧且 $U$ 开, 利用\cref{proposition:紧集的相对紧开邻域}可以选取$\rho>0$,使得
\begin{align*}
\overline{B^{0}(a,r+\rho)}\subseteq U.
\end{align*}
取标准磨光子 $\eta\in C_c^\infty(\mathbb{R}^n)$, 满足
\begin{align*}
\eta\geqslant0,\qquad
\operatorname{spt}\eta\subseteq\overline{B^{0}(0,1)},\qquad
\int_{\mathbb{R}^n}\eta(x)\,\mathrm{d}x=1.
\end{align*}
对 $0<\varepsilon<\rho$, 定义
\begin{align*}
\eta_\varepsilon(x)
\triangleq\frac{1}{\varepsilon^n}
\eta\left(\frac{x}{\varepsilon}\right),\quad 
u_\varepsilon(x)
\triangleq\int_U\eta_\varepsilon(x-y)u(y)\,\mathrm{d}y,
\quad x\in B.
\end{align*}
由\hyperref[theorem:appC-光滑算子的性质]{磨光子的性质\ref{theorem:appC-光滑算子的性质-1}},有
$u_\varepsilon\in C^\infty(B)$, 且
\begin{align}\label{eqegjiojerjieeieieee}
D_i u_\varepsilon(x)
=\int_Uu(y)D_{x_i}\eta_\varepsilon(x-y)\,\mathrm{d}y,
\qquad x\in B.
\end{align}
令$F(x,y)\triangleq\eta_\varepsilon(x-y)$. 由于$\eta_\varepsilon\in C^\infty(\mathbb{R}^n)$, 故$F\in C^\infty(\mathbb{R}^{2n})$.
令$z=x-y$. 因为
\begin{align*}
\partial_{x_i}z_j=\delta_{ij},\quad\partial_{y_i}z_j=-\delta_{ij},
\end{align*}
其中$\delta_{ij}$为 Kronecker 符号, 所以对每个$1\leqslant i\leqslant n$, 多元链式法则给出
\begin{align*}
\partial_{x_i}F(x,y)
&=\sum\limits_{j=1}^{n}(\partial_j\eta_\varepsilon)(x-y)\delta_{ij}
=(\partial_i\eta_\varepsilon)(x-y),\\
\partial_{y_i}F(x,y)
&=-\sum\limits_{j=1}^{n}(\partial_j\eta_\varepsilon)(x-y)\delta_{ij}
=-(\partial_i\eta_\varepsilon)(x-y).
\end{align*}
这里$\partial_i\eta_\varepsilon$表示磨光子对第$i$个自变量的偏导数. 因而
\begin{align}\label{eq::JIFJIOJIOWIJOjjjjjiiiii}
\partial_{x_i}F=-\partial_{y_i}F\iff D_{x_i}\eta_\varepsilon(x-y)
=-D_{y_i}\eta_\varepsilon(x-y).
\end{align}
对于固定的 $x\in B$, 函数
\begin{align*}
y\longmapsto\eta_\varepsilon(x-y)
\end{align*}
光滑, 并且其支集包含于
\begin{align*}
\overline{B^{0}(x,\varepsilon)}
\subseteq\overline{B^{0}(a,r+\rho)}
\subseteq U.
\end{align*}
因此它属于 $C_c^\infty(U)$.再由\eqref{eqwr:Njkojioejioio2}\eqref{eqegjiojerjieeieieee}\eqref{eq::JIFJIOJIOWIJOjjjjjiiiii}式,得
\begin{align*}
D_i u_\varepsilon(x)
=-\int_Uu(y)D_{y_i}\eta_\varepsilon(x-y)\,\mathrm{d}y
=0,\quad \forall x\in B.
\end{align*}
由于 $i$ 任意, 有 $Du_\varepsilon=0$ 在 $B$ 上处处成立.
对于任意 $x,z\in B$, 因为球 $B$ 是凸集, 所以由微积分基本定理得
\begin{align*}
u_{\varepsilon}(x)-u_{\varepsilon}(z)=\int_0^1{\frac{\mathrm{d}}{\mathrm{d}t}u_{\varepsilon}\bigl( z+t(x-z) \bigr)}\,\mathrm{d}t=\int_0^1{Du_{\varepsilon}\bigl( z+t(x-z) \bigr)}\cdot (x-z)\,\mathrm{d}t=0.
\end{align*}
因此存在常数 $c_\varepsilon$, 使得
\begin{align*}
u_\varepsilon(x)=c_\varepsilon,
\quad \forall x\in B.
\end{align*}

由于 $\overline{B}\subseteq U$ 且 $u\in L^1_{\mathrm{loc}}(U)$, 再由\hyperref[theorem:appC-光滑算子的性质]{磨光子的性质\ref{theorem:appC-光滑算子的性质-4}}, 有
\begin{align*}
\lim\limits_{\varepsilon\to0^+}
\|u_\varepsilon-u\|_{L^1(B)}=0.
\end{align*}
定义
\begin{align*}
c_B\triangleq\frac{1}{|B|}\int_Bu(x)\,\mathrm{d}x.
\end{align*}
因为 $u_\varepsilon=c_\varepsilon$ 在 $B$ 上成立, 所以
\begin{align*}
|c_\varepsilon-c_B|
=\left|
\frac{1}{|B|}\int_B
\bigl(u_\varepsilon(x)-u(x)\bigr)\,\mathrm{d}x
\right|
\leqslant
\frac{1}{|B|}
\|u_\varepsilon-u\|_{L^1(B)}.
\end{align*}
从而
\begin{align*}
\|u-c_B\|_{L^1(B)}
\leqslant
\|u-u_\varepsilon\|_{L^1(B)}
+\|u_\varepsilon-c_B\|_{L^1(B)}=\|u-u_\varepsilon\|_{L^1(B)}
+|B||c_\varepsilon-c_B|
\leqslant
2\|u-u_\varepsilon\|_{L^1(B)}.
\end{align*}
令 $\varepsilon\to0^+$, 得
\begin{align*}
\|u-c_B\|_{L^1(B)}=0.
\end{align*}
因此 $u=c_B$ 在 $B$ 上几乎处处成立.

$\mathbf{Step}\mathbf{2.}$下面利用 $U$ 的连通性, 证明这些局部常数一致.
固定一个开球 $B_*$, 满足 $\overline{B_*}\subseteq U$.
由$\mathbf{Step}\mathbf{1}$的结论知, 存在常数 $c$, 使得
\begin{align*}
u=c\text{在 }B_*\text{ 上几乎处处成立}.
\end{align*}
令 $\mathcal{S}$ 为所有满足
\begin{align*}
\overline{B}\subseteq U,\qquad
u=c\text{在 }B\text{ 上几乎处处成立}
\end{align*}
的开球 $B$ 的并集.
由定义, $\mathcal{S}$ 是 $U$ 中的开集, 且
$B_*\subseteq\mathcal{S}$, 所以 $\mathcal{S}\neq\varnothing$.

我们证明 $\mathcal{S}$ 在 $U$ 中也是闭集.
任取 $x\in\overline{\mathcal{S}}\cap U$, 由$U$是开集,可选取开球 $B_x$, 使得
\begin{align*}
x\in B_x,\qquad \overline{B_x}\subseteq U.
\end{align*}
由$\mathbf{Step}\mathbf{1}$的结论, 存在常数$d$, 使得 $u=d$ 在 $B_x$ 上几乎处处成立.
由于 $x\in\overline{\mathcal{S}}$, 而 $B_x$ 是 $x$ 的开邻域, 所以
$B_x\cap\mathcal{S}\neq\varnothing$.
取 $y\in B_x\cap\mathcal{S}$,根据 $\mathcal{S}$的定义知，存在开球 $B_y$, 使得
\begin{align*}
y\in B_y\subseteq \mathcal{S},\quad \overline{B_y}\subseteq U,\quad
u=c\text{在 }B_y\text{ 上几乎处处成立}.
\end{align*}
于是 $O\triangleq B_x\cap B_y$ 是非空开集, 因而 $|O|>0$.
在 $O$ 上, $u=c$ 和 $u=d$ 均几乎处处成立, 所以
\begin{align*}
|c-d||O|=\int_O{|c}-d|\,\mathrm{d}x\leqslant \int_O{|c}-u|\,\mathrm{d}x+\int_O{|u}-d|\,\mathrm{d}x=0.
\end{align*}
故 $c=d$, 从而 $u=c$ 在 $B_x$ 上几乎处处成立.
根据 $\mathcal{S}$ 的定义, 有 $B_x\subseteq\mathcal{S}$, 特别地,
$x\in\mathcal{S}$,即$\overline{\mathcal{S}}\cap U\subseteq \mathcal{S},$
这就证明了 $\mathcal{S}$ 在 $U$ 中闭.

由于 $U$ 连通, 而 $\mathcal{S}$ 非空且在 $U$ 中既开又闭, 若$\mathcal{S}\neq U$,则$U=\mathcal{S} \sqcup (U\setminus \mathcal{S})$,这与连通的定义矛盾!故$\mathcal{S}=U.$
于是, 满足 $\overline{B}\subseteq U$ 且 $u=c$ 在 $B$ 上几乎处处成立的开球构成 $U$ 的一个开覆盖.
由\hyperref[Real Analysis-theorem:Lindelöf定理]{Lindelöf定理}, 可以从中选取可数子覆盖
$\{B_j\}_{j=1}^{\infty}$, 使得
\begin{align*}
U=\bigcup\limits_{j=1}^{\infty}B_j.
\end{align*}
对每个 $j$, 记
\begin{align*}
N_j\triangleq\{x\in B_j:u(x)\neq c\}.
\end{align*}
则 $|N_j|=0$, 并且
\begin{align*}
\{x\in U:u(x)\neq c\}
=\bigcup\limits_{j=1}^{\infty}N_j.
\end{align*}
可数个零测集的并仍为零测集, 因此 $u=c$ 在整个 $U$ 上几乎处处成立.
\end{proof}

\begin{proposition}\label{proposition:W1p空间函数基本性质}
设 $-\infty<a<b<+\infty$, $u\in W^{1,p}(a,b)$, 其中 $1\leqslant p\leqslant\infty$, 并用 $Du$ 表示 $u$ 的弱导数.
\begin{enumerate}[(1)]
\item\label{proposition:W1p空间函数基本性质-1}
存在绝对连续函数 $f\in AC[a,b]$, 使得 $u=f$ 在 $(a,b)$ 上几乎处处成立. 函数 $f$ 的导数 $f'$ 几乎处处存在, 且
\begin{align*}
f'(x)=Du(x),\,\,\text{a.e.}\,x\in (a,b),\quad
f'\in L^p(a,b).
\end{align*}

\item\label{proposition:W1p空间函数基本性质-2}
对几乎处处的 $x,y\in[a,b]$, 有
\begin{align*}
|u(x)-u(y)|
\leqslant
|x-y|^{1-\frac{1}{p}}\|Du(t)\|_{L^p(a,b)}.
\end{align*}
\end{enumerate}
\end{proposition}
\begin{note}
受到\refcor{Real Analysis-corollary:Lebesgue积分基本定理}的启发,可以直接猜测$Du$的变限积分加上某个常数就是所求函数$f$,然后再尝试严格证明.
\end{note}
\begin{proof}
\begin{enumerate}[(1)]
\item \noindent$\mathbf{Step}\mathbf{3}.${\heiti 证明 $Du$ 的变上限积分 $v$ 绝对连续且属于 $L^p(a,b)$.}

由于 $1\leqslant p\leqslant\infty$ 且区间 $(a,b)$ 的测度有限, 由\refpro{Real Analysis-proposition:Lp空间包含关系},有
\begin{align*}
u,Du\in L^p(a,b)\subseteq L^1(a,b).
\end{align*}
定义
\begin{align*}
v(x)\triangleq\int_a^xDu(t)\,\mathrm{d}t,\qquad x\in[a,b].
\end{align*}
由\refcor{Real Analysis-corollary:Lebesgue积分基本定理}, 有 $v\in AC[a,b]$, 且其导数几乎处处存在, 满足
\begin{align*}
v'(x)=Du(x),\quad \mathrm{a.e.}\ x\in(a,b).
\end{align*}
特别地, $v'\in L^p(a,b)$. 此外,
\begin{align*}
|v(x)|
\leqslant\int_a^b|Du(t)|\,\mathrm{d}t<\infty,\qquad x\in[a,b].
\end{align*}
所以 $v$ 有界. 由于 $(a,b)$ 的测度有限, 对全部 $1\leqslant p\leqslant\infty$, 均有 $v\in L^p(a,b)$.

\noindent$\mathbf{Step}\mathbf{2}.${\heiti 证明 $u-v$ 几乎处处为常数.}

任取 $\varphi\in C_c^\infty(a,b)$, 将其在端点 $a,b$ 处取值为零, 仍记为 $\varphi$,则延拓后的函数属于 $C_c^\infty[a,b]$, 且 $\varphi(a)=\varphi(b)=0$. 因为 $v\in AC[a,b]$, 所以 $v\varphi\in AC[a,b]$, 且
\begin{align*}
(v\varphi)'(x)
=v'(x)\varphi(x)+v(x)\varphi'(x),\quad
\mathrm{a.e.}\ x\in(a,b).
\end{align*}
再由\refcor{Real Analysis-corollary:Lebesgue积分基本定理}, 得
\begin{align*}
0=v(b)\varphi(b)-v(a)\varphi(a)
=\int_a^b(v\varphi)'(t)\,\mathrm{d}t=\int_a^b
\left(v'(t)\varphi(t)+v(t)\varphi'(t)\right)\,\mathrm{d}t.
\end{align*}
于是
\begin{align*}
\int_a^bv(t)\varphi'(t)\,\mathrm{d}t
=-\int_a^bv'(t)\varphi(t)\,\mathrm{d}t
=-\int_a^bDu(t)\varphi(t)\,\mathrm{d}t.
\end{align*}
结合 $u$ 的弱导数定义, 有
\begin{align*}
\int_a^b(u-v)\varphi'\,\mathrm{d}t
=\int_a^bu\varphi'\,\mathrm{d}t
-\int_a^bv\varphi'\,\mathrm{d}t
=-\int_a^bDu\,\varphi\,\mathrm{d}t
+\int_a^bDu\,\varphi\,\mathrm{d}t
=0.
\end{align*}
因此 $h\triangleq u-v\in L^p(a,b)\subseteq L^1(a,b)$ 的弱导数在 $(a,b)$ 上为零. 由\refpro{proposition:ch05--examoi-11}, 存在常数 $c\in\mathbb{R}$, 使得
\begin{align*}
u(x)-v(x)=c,\quad \mathrm{a.e.}\ x\in(a,b).
\end{align*}

\noindent$\mathbf{Step}\mathbf{3}.${\heiti 证明 $v+c$ 就是所求的 $f$.}

令
\begin{align*}
f(x)\triangleq v(x)+c,\qquad x\in[a,b].
\end{align*}
则 $f\in AC[a,b]$, 且
\begin{align*}
u(x)=f(x),\quad
f'(x)=v'(x)=Du(x),\qquad
\mathrm{a.e.}\ x\in(a,b).
\end{align*}
因此 $u$ 几乎处处等于绝对连续函数 $f$, 而 $f$ 的导数几乎处处存在, 等于 $u$ 的弱导数, 并且属于 $L^p(a,b)$.

\item 由\ref{proposition:W1p空间函数基本性质-1}的结论, 存在 $f\in AC[a,b]$, 使得
\begin{align*}
f(x)=u(x),\quad
f'(x)=Du(x),\qquad
\mathrm{a.e.}\ x\in(a,b).
\end{align*}
任取 $x,y\in[a,b]$. 当 $x=y$ 时, 所需不等式按上述约定显然成立. 下面先设 $x<y$. 由\hyperref[Real Analysis-theorem:Lebesgue积分基本定理]{Lebesgue积分基本定理}, 得
\begin{align*}
f(y)-f(x)
=\int_x^yf'(t)\,\mathrm{d}t
=\int_x^yDu(t)\,\mathrm{d}t.
\end{align*}
因此
\begin{align*}
|f(y)-f(x)|
\leqslant\int_x^y|Du(t)|\,\mathrm{d}t.
\end{align*}

当$p=1$时, 直接有
\begin{align*}
|f(y)-f(x)|
\leqslant\int_x^y|Du(t)|\,\mathrm{d}t
\leqslant\int_a^b|Du(t)|\,\mathrm{d}t
=\|Du\|_{L^1(a,b)}.
\end{align*}

当 $1<p<\infty$ 时, 令 $q=\frac{p}{p-1}$, 则 $\frac{1}{p}+\frac{1}{q}=1$. 由\hyperref[Real Analysis-theorem:实变函数--Hölder不等式]{Hölder不等式}, 得
\begin{align*}
|f(y)-f(x)|
&\leqslant\int_x^y|Du(t)|\,\mathrm{d}t\leqslant
\left(\int_x^y1^q\,\mathrm{d}t\right)^{\frac{1}{q}}
\left(\int_x^y|Du(t)|^p\,\mathrm{d}t\right)^{\frac{1}{p}}\\
&=(y-x)^{1-\frac{1}{p}}
\left(\int_x^y|Du(t)|^p\,\mathrm{d}t\right)^{\frac{1}{p}}\leqslant
|x-y|^{1-\frac{1}{p}}\|Du\|_{L^p(a,b)}.
\end{align*}

当 $p=\infty$ 时, 由本质上确界的定义,得
\begin{align*}
|Du(t)|\leqslant\|Du\|_{L^\infty(a,b)},\quad
\mathrm{a.e.}\ t\in(a,b).
\end{align*}
因此
\begin{align*}
|f(y)-f(x)|
\leqslant\int_x^y|Du(t)|\,\mathrm{d}t
\leqslant(y-x)\|Du\|_{L^\infty(a,b)}
=|x-y|^{1-\frac{1}{p}}\|Du\|_{L^p(a,b)}.
\end{align*}

当 $y<x$ 时, 交换 $x,y$ 即可. 因此, 对所有 $1\leqslant p\leqslant\infty$ 及所有 $x,y\in[a,b]$, 均有
\begin{align*}
|f(x)-f(y)|
\leqslant
|x-y|^{1-\frac{1}{p}}\|Du\|_{L^p(a,b)},
\end{align*}

最后, 记
\begin{align*}
N\triangleq
\{x\in(a,b):u(x)\neq f(x)\}\cup\{a,b\}.
\end{align*}
因为 $u=f$ 在 $(a,b)$ 上几乎处处成立, 所以 $|N|=0$. 对于任意 $x,y\in[a,b]\setminus N$, 有
\begin{align*}
|u(x)-u(y)|
=|f(x)-f(y)|
\leqslant
|x-y|^{1-\frac{1}{p}}\|Du\|_{L^p(a,b)}.
\end{align*}
\end{enumerate}
\end{proof}

\begin{theorem}[Sobolev空间都是Banach空间]\label{theorem:ch05-WkpBanach}
设 $U \subseteq \mathbb{R}^n$ 是一个非空开集,对每个 $k=1,2,\cdots$ 及 $1\leqslant p\leqslant\infty$，Sobolev空间 $W^{k,p}(U)$ 是Banach空间。
\end{theorem}
\begin{proof}

\noindent$\mathbf{Step}\mathbf{1}.${\heiti 首先验证 $\|u\|_{W^{k,p}(U)}$ 是范数。}

(i)齐次性和正定性:显然
\begin{align*}
\|\lambda u\|_{W^{k,p}(U)}=|\lambda|\,\|u\|_{W^{k,p}(U)},
\end{align*}
且\begin{align*}
\|u\|_{W^{k,p}(U)}=0 \Longleftrightarrow \|D^{\alpha}u\|_{L^p(U)}=0,\quad |\alpha|\leqslant k.
\end{align*}
特别地, 取 $\alpha=(0,\cdots,0)$, 有
\begin{align*}
\|u\|_{L^p(U)}=0 \Longleftrightarrow u(x)=0\ \mathrm{a.e.}\ x\in U.
\end{align*}
又当 $u(x)=0\ \mathrm{a.e.}\ x\in U$ 时, 显然有
\begin{align*}
\|D^{\alpha}u\|_{L^p(U)}=0,\quad |\alpha|\leqslant k \Longrightarrow \|u\|_{W^{k,p}(U)}=0.
\end{align*}
故
\begin{align*}
\|u\|_{W^{k,p}(U)}=0 \Longleftrightarrow u(x)=0\ \mathrm{a.e.}\ x\in U.
\end{align*}

(ii)三角不等式:现设 $u,v\in W^{k,p}(U)$。则当 $1\leqslant p<\infty$ 时，\hyperref[Real Analysis-theorem:实变函数--Minkowski不等式]{Minkowski不等式}和\hyperref[Real Analysis-theorem:离散Minkowski不等式]{离散Minkowski不等式}给出
\begin{align*}
\|u+v\|_{W^{k,p}(U)}
&=\left(\sum_{|\alpha|\leqslant k}\|D^{\alpha}u+D^{\alpha}v\|_{L^{p}(U)}^{p}\right)^{\frac{1}{p}}\leqslant\left(\sum_{|\alpha|\leqslant k}\left(\|D^{\alpha}u\|_{L^{p}(U)}+\|D^{\alpha}v\|_{L^{p}(U)}\right)^{p}\right)^{\frac{1}{p}}\\
&\leqslant\left(\sum_{|\alpha|\leqslant k}\|D^{\alpha}u\|_{L^{p}(U)}^{p}\right)^{\frac{1}{p}}
+\left(\sum_{|\alpha|\leqslant k}\|D^{\alpha}v\|_{L^{p}(U)}^{p}\right)^{\frac{1}{p}}=\|u\|_{W^{k,p}(U)}+\|v\|_{W^{k,p}(U)}.
\end{align*}
故 \(1\leqslant p<\infty\) 时的三角不等式成立。

当 \(p=\infty\) 时，根据定义有
\[
\|u\|_{W^{k,\infty}(U)}
\triangleq \sum_{|\alpha|\leqslant k}\mathrm{ess}\,\mathrm{sup}\,_U |D^{\alpha}u|.
\]
先证明 \(L^{\infty}\) 的本质上确界满足三角不等式。设 \(f,g\in L^{\infty}(U)\)，记
\[
M_f\triangleq \mathrm{ess}\,\mathrm{sup}\,_U |f|,\qquad M_g\triangleq \mathrm{ess}\,\mathrm{sup}\,_U |g|.
\]
由本质上确界的定义，\(|f(x)|\leqslant M_f\) 且 \(|g(x)|\leqslant M_g\) 对几乎处处的 \(x\in U\) 成立。于是
\[
|f(x)+g(x)|\leqslant |f(x)|+|g(x)|\leqslant M_f+M_g
\]
对几乎处处的 \(x\in U\) 成立。故 \(M_f+M_g\) 是 \(|f+g|\) 的一个本质上界，从而
\[
\mathrm{ess}\,\mathrm{sup}\,_U |f+g|\leqslant M_f+M_g
=\mathrm{ess}\,\mathrm{sup}\,_U |f|+\mathrm{ess}\,\mathrm{sup}\,_U |g|.
\]

现设\(u,v\in W^{k,\infty}(U)\)。对每个多重指标 \(|\alpha|\leqslant k\)，弱导数满足
\[
D^{\alpha}(u+v)=D^{\alpha}u+D^{\alpha}v.
\]
由上述 \(L^{\infty}\) 本质上确界的三角不等式，有
\[
\mathrm{ess}\,\mathrm{sup}\,_U |D^{\alpha}(u+v)|
\leqslant \mathrm{ess}\,\mathrm{sup}\,_U |D^{\alpha}u|
+\mathrm{ess}\,\mathrm{sup}\,_U |D^{\alpha}v|.
\]
对所有 \(|\alpha|\leqslant k\) 求和，得到
\begin{align*}
\|u+v\|_{W^{k,\infty}(U)}
&=\sum_{|\alpha|\leqslant k}\mathrm{ess}\,\mathrm{sup}\,_U |D^{\alpha}(u+v)|\\
&\leqslant \sum_{|\alpha|\leqslant k}\left(\mathrm{ess}\,\mathrm{sup}\,_U |D^{\alpha}u|+\mathrm{ess}\,\mathrm{sup}\,_U |D^{\alpha}v|\right)\\
&=\sum_{|\alpha|\leqslant k}\mathrm{ess}\,\mathrm{sup}\,_U |D^{\alpha}u|
+\sum_{|\alpha|\leqslant k}\mathrm{ess}\,\mathrm{sup}\,_U |D^{\alpha}v|\\
&=\|u\|_{W^{k,\infty}(U)}+\|v\|_{W^{k,\infty}(U)}.
\end{align*}
故 \(p=\infty\) 时的三角不等式成立。

\noindent$\mathbf{Step}\mathbf{2}.${\heiti 证明 $W^{k,p}(U)$ 完备。}

设 $\{u_{m}\}_{m=1}^{\infty}$ 为 $W^{k,p}(U)$ 中的 Cauchy 序列。则对每个 $|\alpha|\leqslant k$，$\{D^{\alpha}u_{m}\}_{m=1}^{\infty}$ 是 $L^{p}(U)$ 中的 Cauchy 序列。由于 $L^{p}(U)$ 完备，存在函数 $u_{\alpha}\in L^{p}(U)$ 使对每个 $|\alpha|\leqslant k$，在$L^{p}(U)$中,都有
\begin{align*}
D^{\alpha}u_{m}\to u_{\alpha}.
\end{align*}
特别地，在$L^{p}(U)$中,有$u_{m}\to u_{(0,\cdots,0)}\triangleq u$.

我们现在断言
\begin{align}
u\in W^{k,p}(U),\qquad D^{\alpha}u=u_{\alpha}\quad (|\alpha|\leqslant k).
\label{align:ch05-完备性结论}
\end{align}
为验证此断言，先证明两个极限等式.
设 $p'$ 为 $p$ 的共轭指数。由\hyperref[Real Analysis-theorem:实变函数--Hölder不等式]{Hölder不等式}，对任意 $\phi\in C_{c}^{\infty}(U)$，有
\begin{align*}
\left|\int_{U}(u_{m}-u)D^{\alpha}\phi\,\mathrm{d}x\right|
&\leqslant\|u_{m}-u\|_{L^{p}(U)}\|D^{\alpha}\phi\|_{L^{p'}(U)},\\
\left|\int_{U}(D^{\alpha}u_{m}-u_{\alpha})\phi\,\mathrm{d}x\right|
&\leqslant\|D^{\alpha}u_{m}-u_{\alpha}\|_{L^{p}(U)}\|\phi\|_{L^{p'}(U)}.
\end{align*}
因为 $\phi\in C_{c}^{\infty}(U)$，所以 $\|D^{\alpha}\phi\|_{L^{p'}(U)}$ 与 $\|\phi\|_{L^{p'}(U)}$ 均为有限常数。令 $m\to\infty$，由$u_m\to u$ 且 $D^{\alpha}u_m\to u_{\alpha}$ 于 $L^{p}(U)$可知两式右端均趋于 $0$，于是
\begin{align*}
\lim_{m\to\infty}\int_{U}u_{m}D^{\alpha}\phi\,\mathrm{d}x=\int_{U}uD^{\alpha}\phi\,\mathrm{d}x,\qquad
\lim_{m\to\infty}\int_{U}D^{\alpha}u_{m}\phi\,\mathrm{d}x=\int_{U}u_{\alpha}\phi\,\mathrm{d}x.
\end{align*}
再结合弱导数的定义得
\begin{align*}
\int_{U}uD^{\alpha}\phi\,\mathrm{d}x
=\lim_{m\to\infty}\int_{U}u_{m}D^{\alpha}\phi\,\mathrm{d}x=\lim_{m\to\infty}(-1)^{|\alpha|}\int_{U}D^{\alpha}u_{m}\phi\,\mathrm{d}x=(-1)^{|\alpha|}\int_{U}u_{\alpha}\phi\,\mathrm{d}x.
\end{align*}
故\eqref{align:ch05-完备性结论}式成立。由于对一切 $|\alpha|\leqslant k$ 都有 $D^{\alpha}u_{m}\to D^{\alpha}u$ 于 $L^{p}(U)$，可见 $u_{m}\to u$ 于 $W^{k,p}(U)$，即得所需。
\end{proof}

\begin{proposition}[$H^k(U)$ 都是Hilbert空间]\label{proposition:ch05-Hk-Hilbert-a7d3}
设 $U\subseteq\mathbb{R}^n$ 为非空开集, $k$ 为非负整数.  则 $H^k(U)\triangleq W^{k,2}(U)$ 关于内积
\begin{align*}
\langle u,v\rangle_{H^k(U)}
\triangleq
\sum\limits_{|\alpha|\leqslant k}
\int_U D^\alpha u\,D^\alpha v\,\mathrm{d}x
\end{align*}
构成实 Hilbert 空间, 其中约定 $D^0u=u$. 此内积诱导的范数恰为
\begin{align*}
\|u\|_{H^k(U)}
=
\left(
\sum\limits_{|\alpha|\leqslant k}
\|D^\alpha u\|_{L^2(U)}^2
\right)^{\frac{1}{2}}.
\end{align*}
此外, $H_0^k(U)$ 关于上述内积也是 Hilbert 空间.
\end{proposition}
\begin{proof}

(1)证明$H^k(U)$是Hilbert 空间.

\noindent$\mathbf{Step}\mathbf{1}.${\heiti 证明$\langle \cdot,\cdot\rangle_{H^k(U)}$是内积.}

由弱导数的线性性, $H^k(U)$ 为实向量空间. 对任意 $u,v\in H^k(U)$ 及 $|\alpha|\leqslant k$, 由 Cauchy--Schwarz 不等式得
\begin{align*}
\left|
\int_U D^\alpha u\,D^\alpha v\,\mathrm{d}x
\right|
\leqslant
\|D^\alpha u\|_{L^2(U)}
\|D^\alpha v\|_{L^2(U)}
<\infty.
\end{align*}
因此 $\langle u,v\rangle_{H^k(U)}$是良定义的, 且显然满足双线性和对称性. 又有
\begin{align*}
\langle u,u\rangle_{H^k(U)}
=
\sum\limits_{|\alpha|\leqslant k}
\|D^\alpha u\|_{L^2(U)}^2
\geqslant
\|u\|_{L^2(U)}^2
\geqslant 0.
\end{align*}
若 $\langle u,u\rangle_{H^k(U)}=0$, 则 $\|u\|_{L^2(U)}=0$, 即 $u=0$ 几乎处处成立. 反之, 若 $u=0$ 几乎处处成立, 则其各阶弱导数均为零, 从而 $\langle u,u\rangle_{H^k(U)}=0$. 故上述表达式确为内积, 且其诱导范数就是所给的 Sobolev 范数.

\noindent$\mathbf{Step}\mathbf{2}.${\heiti 证明$H^k(U)$的完备性.}

设 $\{u_m\}_{m=1}^{\infty}$ 为 $H^k(U)$ 中的 Cauchy 序列. 对每个满足 $|\alpha|\leqslant k$ 的多重指标 $\alpha$, 有
\begin{align*}
\|D^\alpha u_m-D^\alpha u_\ell\|_{L^2(U)}
\leqslant
\|u_m-u_\ell\|_{H^k(U)}.
\end{align*}
因此 $\{D^\alpha u_m\}_{m=1}^{\infty}$ 为 $L^2(U)$ 中的 Cauchy 序列. 由 $L^2(U)$ 的完备性, 存在 $v_\alpha\in L^2(U)$, 使得
\begin{align*}
\lim\limits_{m\to\infty}
\|D^\alpha u_m-v_\alpha\|_{L^2(U)}=0.
\end{align*}
特别地, 令 $u\triangleq v_0$, 则 $u_m\to u$ 于 $L^2(U)$, 且$u\in L^2(U)$.

下面验证 $v_\alpha$ 恰为 $u$ 的 $\alpha$ 阶弱导数. 任取 $\varphi\in C_c^\infty(U)$. 由 $u_m$ 的弱导数定义得
\begin{align}\label{eqHUHIUHIUHUHUHIOJIO}
\int_U u_mD^\alpha\varphi\,\mathrm{d}x
=
(-1)^{|\alpha|}
\int_U D^\alpha u_m\,\varphi\,\mathrm{d}x.
\end{align}
由$\varphi\in C_c^\infty(U)$知$\varphi,D^\alpha\varphi\in L^2(U)$, 再由Cauchy--Schwarz 不等式得
\begin{align*}
\left|
\int_U (u_m-u)D^\alpha\varphi\,\mathrm{d}x
\right|
&\leqslant
\|u_m-u\|_{L^2(U)}
\|D^\alpha\varphi\|_{L^2(U)}
\longrightarrow 0,\\
\left|
\int_U (D^\alpha u_m-v_\alpha)\varphi\,\mathrm{d}x
\right|
&\leqslant
\|D^\alpha u_m-v_\alpha\|_{L^2(U)}
\|\varphi\|_{L^2(U)}
\longrightarrow 0.
\end{align*}
故对\eqref{eqHUHIUHIUHUHUHIOJIO}式令$m\to\infty$, 得
\begin{align*}
\int_U uD^\alpha\varphi\,\mathrm{d}x
=
(-1)^{|\alpha|}
\int_U v_\alpha\varphi\,\mathrm{d}x.
\end{align*}
这说明 $D^\alpha u=v_\alpha$ 在弱意义下成立. 因此 $u\in H^k(U)$.

由于满足 $|\alpha|\leqslant k$ 的多重指标只有有限个, 故
\begin{align*}
\|u_m-u\|_{H^k(U)}^2
=
\sum\limits_{|\alpha|\leqslant k}
\|D^\alpha u_m-D^\alpha u\|_{L^2(U)}^2
=
\sum\limits_{|\alpha|\leqslant k}
\|D^\alpha u_m-v_\alpha\|_{L^2(U)}^2
\longrightarrow 0.
\end{align*}
因此 $H^k(U)$ 完备, 从而为 Hilbert 空间.

(2)现在证明$H_0^k(U)$也为 Hilbert 空间.

\noindent$\mathbf{Step}\mathbf{1}.${\heiti 验证 $C_c^\infty(U)$ 是 $H^k(U)$ 的线性子空间.}

任取 $\varphi\in C_c^\infty(U)$, 记 $K\triangleq\operatorname{spt}\varphi$. 则 $K$ 是包含于 $U$ 的紧集, 因而 $m(K)<\infty$.

对每个满足 $|\alpha|\leqslant k$ 的多重指标 $\alpha$, 函数 $D^\alpha\varphi$ 连续, 且
\begin{align*}
\operatorname{spt}(D^\alpha\varphi)\subseteq K.
\end{align*}
因此 $D^\alpha\varphi$ 有界, 可取常数 $M_\alpha<\infty$, 使得
\begin{align*}
|D^\alpha\varphi(x)|\leqslant M_\alpha,
\quad x\in U.
\end{align*}
于是
\begin{align*}
\int_U |D^\alpha\varphi|^2\,\mathrm{d}x
=
\int_K |D^\alpha\varphi|^2\,\mathrm{d}x
\leqslant M_\alpha^2m(K)
<\infty.
\end{align*}
此外, 对任意 $\eta\in C_c^\infty(U)$, 反复分部积分得到
\begin{align*}
\int_U \varphi D^\alpha\eta\,\mathrm{d}x
=
(-1)^{|\alpha|}
\int_U D^\alpha\varphi\,\eta\,\mathrm{d}x.
\end{align*}
因此 $\varphi$ 的常义导数也是相应的弱导数, 故 $\varphi\in H^k(U)$. 这证明了
\begin{align*}
C_c^\infty(U)\subseteq H^k(U).
\end{align*}

任取 $\varphi,\psi\in C_c^\infty(U)$ 及 $a,b\in\mathbb{R}$. 函数 $a\varphi+b\psi$ 光滑, 且
\begin{align*}
\operatorname{spt}(a\varphi+b\psi)
\subseteq
\operatorname{spt}\varphi\cup\operatorname{spt}\psi.
\end{align*}
右端是包含于 $U$ 的紧集, 故 $a\varphi+b\psi\in C_c^\infty(U)$. 再结合 $0\in C_c^\infty(U)$, 可知 $C_c^\infty(U)$ 是 $H^k(U)$ 的线性子空间.

\noindent$\mathbf{Step}\mathbf{2}.${\heiti 验证闭包 $H_0^k(U)$ 仍是线性子空间.}

任取 $u,v\in H_0^k(U)$ 及 $a,b\in\mathbb{R}$. 根据闭包的定义, 存在 $\varphi_m,\psi_m\in C_c^\infty(U)$, 使得
\begin{align*}
\|\varphi_m-u\|_{H^k(U)}\longrightarrow 0,
\quad
\|\psi_m-v\|_{H^k(U)}\longrightarrow 0.
\end{align*}
由上一步, $a\varphi_m+b\psi_m\in C_c^\infty(U)$, 且
\begin{align*}
\|a\varphi_m+b\psi_m-(au+bv)\|_{H^k(U)}
\leqslant
|a|\|\varphi_m-u\|_{H^k(U)}
+
|b|\|\psi_m-v\|_{H^k(U)}
\longrightarrow 0.
\end{align*}
因此 $au+bv$ 也属于 $C_c^\infty(U)$ 在 $H^k(U)$ 中的闭包, 即 $au+bv\in H_0^k(U)$. 又 $0\in H_0^k(U)$, 故 $H_0^k(U)$ 是线性子空间.

验证 $H_0^k(U)$ 在 $H^k(U)$ 中是闭集.
设 $u_m\in H_0^k(U)$, 且存在 $u\in H^k(U)$, 使得
\begin{align*}
\|u_m-u\|_{H^k(U)}\longrightarrow 0.
\end{align*}
由于每个 $u_m$ 都属于 $C_c^\infty(U)$ 的闭包, 对每个正整数 $m$, 可以选取 $\varphi_m\in C_c^\infty(U)$, 使得
\begin{align*}
\|\varphi_m-u_m\|_{H^k(U)}<\frac{1}{m}.
\end{align*}
于是
\begin{align*}
\|\varphi_m-u\|_{H^k(U)}
\leqslant
\|\varphi_m-u_m\|_{H^k(U)}
+
\|u_m-u\|_{H^k(U)}
<
\frac{1}{m}+\|u_m-u\|_{H^k(U)}
\longrightarrow 0.
\end{align*}
因此 $u$ 仍属于 $C_c^\infty(U)$ 的闭包, 即 $u\in H_0^k(U)$.
这说明 $\overline{H_0^k(U)}\subseteq H_0^k(U)$,这等价于 $H_0^k(U)$ 为闭集.

\noindent$\mathbf{Step}\mathbf{3}.${\heiti 验证 $H_0^k(U)$ 关于继承的内积完备.}

将 $H^k(U)$ 的内积限制在其线性子空间 $H_0^k(U)$ 上, 双线性、对称性和正定性仍然成立, 因而得到 $H_0^k(U)$ 上的内积. 该内积诱导的范数满足
\begin{align*}
\|u\|_{H_0^k(U)}
=
\sqrt{\langle u,u\rangle_{H_0^k(U)}}
=
\|u\|_{H^k(U)},
\quad u\in H_0^k(U).
\end{align*}

任取 $H_0^k(U)$ 中的 Cauchy 序列 $\{u_m\}_{m=1}^{\infty}$. 由于两个空间中的范数相同, 它也是 $H^k(U)$ 中的 Cauchy 序列. 根据 $H^k(U)$ 的完备性, 存在 $u\in H^k(U)$, 使得
\begin{align*}
\|u_m-u\|_{H^k(U)}\longrightarrow 0.
\end{align*}
由上一步证明的$H_0^k(U)$是$H^k(U)$中的闭集, 有 $u\in H_0^k(U)$. 因而
\begin{align*}
\|u_m-u\|_{H_0^k(U)}
=
\|u_m-u\|_{H^k(U)}
\longrightarrow 0.
\end{align*}
因此 $H_0^k(U)$ 中的每个 Cauchy 序列都收敛于该空间中的元素, 即 $H_0^k(U)$ 完备. 故 $H_0^k(U)$ 为 Hilbert 空间.
\end{proof}













\end{document}
